%% 
\documentclass[acmsmall]{acmart}
% Add this to override the default font size
\makeatletter
\renewcommand\normalsize{%
   \@setfontsize\normalsize{9}{11pt} % Adjust font size and line spacing
   \abovedisplayskip 6pt plus 2pt minus 2pt
   \belowdisplayskip \abovedisplayskip
   \abovedisplayshortskip 0pt plus 2pt
   \belowdisplayshortskip 3pt plus 2pt minus 2pt
}
\makeatother
%%

% ============================================================
\usepackage{float} 
\usepackage{array}
\usepackage{colortbl}
\usepackage{xcolor}
\usepackage{graphicx}

\let\Bbbk\relax
\usepackage{amssymb}
\usepackage{pifont}
\usepackage{tikz}
\usepackage[table]{xcolor}
\usepackage{pdflscape}
\usepackage{longtable}
\definecolor{gmrule}{gray}{0.65}
\usetikzlibrary{arrows.meta}

%\usepackage{pdflscape}
\usepackage{longtable} 
\usepackage{rotating}
\usepackage{float}
%\usepackage{amssymb}
\usepackage{xcolor}
\usepackage{tabularx}
\usepackage{makecell}
\usepackage{geometry}
\usepackage{float}
\usepackage{multirow}
\usepackage{graphicx}
\usepackage{subcaption}
\usepackage{hyperref}
\usepackage{lscape}
\usepackage{rotating}
\usepackage{graphicx}
\usepackage{adjustbox}
\usepackage[table]{xcolor} % Add this to the preamble for color support
%%
\usepackage{booktabs}
\usepackage{array}
\usepackage{makecell}
\usepackage[table]{xcolor}
\usepackage{pifont}

% ============================================================
% [REV] Clean version: all revision colouring removed.
% Every change carries a LaTeX comment starting with [REV-Rxx];
% replaced original text is kept as comments starting with [REV-Rxx-ORIGINAL].
% ============================================================
\newcommand{\revadd}[1]{#1}
\newenvironment{revblock}{\par}{\par}

\newcommand{\cmark}{\ding{51}}
\newcommand{\xmark}{\ding{55}}
\newcommand{\pmark}{$\triangle$}

\usepackage{rotating}
\usepackage{multirow}
\usepackage{booktabs}
\usepackage{pifont}

%% Submission ID.
%% Use this when submitting an article to a sponsored event. You'll
%% receive a unique submission ID from the organizers
%% of the event, and this ID should be used as the parameter to this command.
%%\acmSubmissionID{123-A56-BU3}
\acmJournal{CSUR}

%%
%% For managing citations, it is recommended to use bibliography
%% files in BibTeX format.
%%
%% You can then either use BibTeX with the ACM-Reference-Format style,
%% or BibLaTeX with the acmnumeric or acmauthoryear sytles, that include
%% support for advanced citation of software artefact from the
%% biblatex-software package, also separately available on CTAN.
%%
%% Look at the sample-*-biblatex.tex files for templates showcasing
%% the biblatex styles.
%%

%%
%% The majority of ACM publications use numbered citations and
%% references.  The command \citestyle{authoryear} switches to the
%% "author year" style.
%%red}
%% If you are preparing content for an event
%% sponsored by ACM SIGGRAPH, you must use the "author year" style of
%% citations and references.
%% Uncommenting
%% the next command will enable that style.
%%\citestyle{acmauthoryear}


%%
%% end of the preamble, start of the body of the document source.
\begin{document}

%%
%% The "title" command has an optional parameter,
%% allowing the author to define a "short title" to be used in page headers.
\title{Exploring Graph Mamba: A Comprehensive Survey on State-Space Models for Graph Learning}

%%
%% The "author" command and its associated commands are used to define
%% the authors and their affiliations.
%% Of note is the shared affiliation of the first two authors, and the
%% "authornote" and "authornotemark" commands
%% used to denote shared contribution to the research.
 

% Safa Ben Atitallah
\author{Safa Ben Atitallah}
\orcid{0000-0003-0796-3507}
\email{satitallah@psu.edu.sa}
\affiliation{%
  \institution{Robotics \& Internet of Things Lab,
  Research and Initiatives Center, Prince Sultan University}
  \city{Riyadh}
  \country{Saudi Arabia}
}
\affiliation{%
  \institution{University of Manouba}
  \city{Manouba}
  \country{Tunisia}
}

% Chaima Ben Rabah
\author{Chaima Ben Rabah}
\affiliation{%
  \institution{Weill Cornell Medicine}
  \city{Doha}
  \country{Qatar}
}
\affiliation{%
  \institution{University of Manouba}
  \city{Manouba}
  \country{Tunisia}
}

% Maha Driss
\author{Maha Driss}
\orcid{0000-0001-8236-8746}
\affiliation{%
  \institution{Department of Software Engineering,
  College of Engineering and Advanced Computing,
  Alfaisal University}
  \city{Riyadh}
  \country{Saudi Arabia}
}
\affiliation{%
  \institution{University of Manouba}
  \city{Manouba}
  \country{Tunisia}
}

% Wadii Boulila
\author{Wadii Boulila}
\orcid{0000-0003-2133-0757}
\affiliation{%
  \institution{Robotics \& Internet of Things Lab,
  Research and Initiatives Center, Prince Sultan University}
  \city{Riyadh}
  \country{Saudi Arabia}
}
\affiliation{%
  \institution{University of Manouba}
  \city{Manouba}
  \country{Tunisia}
}

% Anis Koubaa
\author{Anis Koubaa}
\orcid{0000-0003-3787-7423}
\affiliation{%
  \institution{Department of Software Engineering,
  College of Engineering and Advanced Computing,
  Alfaisal University}
  \city{Riyadh}
  \country{Saudi Arabia}
}

\renewcommand{\shortauthors}{Ben Atitallah et al.}

%%
%% The abstract is a short summary of the work to be presented in the
%% article.
\begin{abstract}
% [REV-R01] One sentence added (red) announcing the four design questions that organize the survey.
Recently, Graph Mamba has emerged as a promising approach for extending selective State Space Models (SSMs) to graph-structured data, with the aim of efficiently capturing long-range dependencies in graphs. This survey provides a systematic review of Graph Mamba, with particular emphasis on its architectural developments, applications, comparative performance, limitations, and future research directions. We first introduce the foundations of SSMs and Mamba necessary to understand their adaptation to graph-structured data. Then, we examine representative Graph Mamba architectures and focus on their graph-to-sequence construction, ordering strategies, structural encoding, and graph-aware state-space mechanisms. {We organize these architectures around four design questions: what the state-space model reads (tokenization), in what order it reads it (ordering), how deeply graph structure shapes the state dynamics (coupling), and how the state-space model works alongside graph propagation (integration).} We also review recent applications of Graph Mamba across different domains and graph learning tasks. Furthermore, we provide benchmark and comparative analyses to examine the performance and characteristics of existing Graph Mamba approaches across different tasks. Finally, we discuss current limitations, open challenges, and emerging research directions, including opportunities for improving scalability, structural modeling, interpretability, and applicability. This survey provides an architectural and comparative perspective on the rapidly evolving field of Graph Mamba and aims to support further research on selective state-space modeling for graph-structured data.
\end{abstract}
 
%%
%% Keywords. The author(s) should pick words that accurately describe
%% the work being presented. Separate the keywords with commas.
\keywords{State Space Models,
Selective State Space Models,
Mamba,
Graph Mamba,
Graph Learning,
Graph Serialization }

%\received{20 February 2007}
%\received[revised]{12 March 2009}
%\received[accepted]{5 June 2009}

%%
%% This command processes the author and affiliation and title
%% information and builds the first part of the formatted document.
\maketitle

\section{Introduction}

Graph-structured data arises naturally in many real-world systems, including
social networks, molecular structures, and citation networks
\cite{xia2021graph,jiang2022graph,wang2024graph}.
% [REV-NOTE] jiang2022graph (traffic survey) and wang2024graph (Graph-Mamba) do not directly support this generic sentence; consider general graph-learning references. 
In these systems, entities
and their interactions are represented as nodes and edges, enabling complex
relational information to be modeled explicitly. With the advancement of Deep
Learning (DL), Graph Neural Networks (GNNs) have emerged as a widely used
approach for learning from graph-structured data \cite{xia2021graph}. Through
message passing, GNNs aggregate information from neighboring nodes to learn
structural and relational representations, enabling end-to-end learning across
a wide range of graph learning tasks and application domains
\cite{wang2024graph}.

However, GNNs face several limitations, particularly in capturing long-range
dependencies and scaling to large graphs. Their reliance on local message
passing requires information to propagate across multiple layers, which may
increase computational and memory costs \cite{liu2024federated}.
% [REV-NOTE] \cite{liu2024federated} is a federated-GNN survey; consider a reference on over-squashing / long-range limits of message passing instead. 
These
challenges become more pronounced in large and dynamic graphs, where repeated
neighborhood aggregation and frequent representation updates can be
computationally demanding. Graph Transformers address long-range interactions
through global attention, but this can introduce quadratic complexity with
respect to the number of nodes. Architectures such as GraphGPS
\cite{rampavsek2022recipe} and Exphormer \cite{shirzad2023exphormer} combine
local graph processing with global attention-based mechanisms, providing an
important architectural basis for subsequent Graph Mamba\footnote{Throughout
this survey, \emph{Graph Mamba} denotes the family of models that adapt Mamba
and selective state-space models to graph-structured data. \emph{Graph-Mamba}
refers to the specific model of Wang et al.~\cite{wang2024graph}, and
\emph{GMN} refers to the Graph Mamba Network of Behrouz and
Hashemi~\cite{behrouz2024graph}. Three unrelated models published under the
name \emph{GraphMamba} are distinguished as GraphMamba
(HSI)~\cite{yang2024graphmamba} for hyperspectral image classification,
GraphMamba (WSI)~\cite{zheng2025graphmamba} for whole slide image
classification, and GraphMamba (tok.)~\cite{ahmad2025graphmamba} for
graph-guided spectral--spatial tokenization.} models. In this
context, Graph Mamba can be viewed not only as an alternative to conventional
message-passing GNNs, but also as an approach for replacing or complementing
global attention with selective state-space processing for efficient
long-range graph modeling.

% [REV-NOTE] "Originally introduced for modeling dynamical systems" is supported by S4 only indirectly; consider adding a classical reference (Kalman, 1960).
Addressing these limitations has motivated the exploration of State-Space
Models (SSMs) \cite{gu2021efficiently} for graph learning. Originally
introduced for modeling dynamical systems, SSMs have been increasingly adopted
in Machine Learning (ML) due to their ability to efficiently capture long-range
dependencies in sequential data. Their structured state representation and
efficient information propagation provide a suitable foundation for extending
state-space modeling to graph-structured data.

Building on these developments, Mamba \cite{gu2023mamba} has emerged as a
selective SSM architecture for efficient long-sequence modeling. Unlike
Transformers, where self-attention scales quadratically with sequence length,
Mamba employs selective state-space processing whose computational cost scales
linearly with sequence length. Its
input-dependent state-space parameters allow information propagation to be
dynamically controlled according to the input, providing an efficient mechanism
for modeling long-range dependencies.

Graph Mamba extends Mamba and selective SSM mechanisms to graph-structured
data. Since Mamba is inherently designed for ordered sequences, its application
to graphs requires mechanisms that account for graph topology and the absence
of a natural node ordering. Existing Graph Mamba architectures address this
challenge through strategies such as graph-aware serialization, structural
encoding, node ordering, bidirectional scanning, and the integration of local
graph processing with global state-space modeling. These mechanisms aim to
combine the structural modeling capabilities of graph-based methods with the
efficient long-range dependency modeling of selective SSMs.

The rapid growth of Graph Mamba has produced a diverse set of architectures
that differ substantially in how graph structure is incorporated into
state-space computation. Some methods embed topology directly into the
state-space operator, whereas others construct graph-derived tokens, impose
structure-aware serialization, condition state transitions on graph
information, or combine selective SSMs with local message passing. These
differences affect permutation behavior, long-range information propagation,
computational cost, and suitability for different graph settings. Accordingly,
this survey goes beyond cataloguing individual models and organizes the field
around the mechanisms through which graph structure interacts with selective
state-space processing.
% [REV-R02] Added (red): the four design questions are announced here and used as the organizing lens of the survey.
We show that these mechanisms can be analyzed through four architectural design questions: what the state-space model reads (tokenization), in what order it reads it (ordering), how deeply graph structure shapes the state dynamics (coupling), and how the state-space model interacts with graph propagation (integration). These questions provide the architectural framework used throughout the survey to organize the taxonomy and support the analysis of applications, empirical evidence, and open research challenges.

%\begin{figure}[ht]
%    \centering
%    \includegraphics[width=8 cm]{Figures/Picture1.png}
%    \caption{Publication Trend of Graph Mamba.}
%    \label{fig1}
%\end{figure} 


\subsection{Related Surveys and Scope}

The survey literature most relevant to this work encompasses four interconnected
research directions: GNNs, Graph Transformers, SSMs/Mamba, and Graph Mamba. 
Several surveys provide comprehensive perspectives on GNN architectures, applications, and practical challenges. Wu et al. \cite{wu2020comprehensive} established a broad taxonomy covering recurrent, convolutional, graph autoencoder, and spatiotemporal GNNs, while subsequent reviews examined recent architectural developments and applications \cite{khemani2024review}. More recent surveys have focused on practical challenges. Ju et al. \cite{ju2025survey} investigated GNN reliability and robustness under data imbalance, noise, privacy constraints, and out-of-distribution scenarios. Zhang et al. \cite{zhang2026survey} reviewed GNN acceleration from algorithmic, system-level, and customized hardware perspectives. These surveys cover GNN foundations, applications, robustness, scalability, and computational efficiency. However, they do not specifically examine the integration of selective state-space mechanisms with graph-structured data.

Graph Transformers represent another closely related research direction,
particularly for modeling long-range dependencies in graphs. Shehzad et
al.~\cite{shehzad2026graph} provided a comprehensive survey of Graph
Transformers, examining how graph inductive biases and graph attention
mechanisms are integrated into Transformer architectures. Their survey also
classifies existing approaches according to depth, scalability, and
pre-training strategies and discusses applications across node-, edge-, and
graph-level tasks. 
While Graph Transformers primarily rely on attention mechanisms for global
information exchange, Graph Mamba explores selective state-space processing as
an alternative mechanism for efficient long-range graph modeling.


A parallel body of surveys has examined SSMs and Mamba from broader architectural and application perspectives. Patro et al. \cite{patro2024mamba}, Qu et al. \cite{qu2024survey}, and Wang et al. \cite{wang2024state} review the foundations and evolution of structured SSMs and Mamba, highlighting their efficiency in long-sequence modeling and their applications across multiple domains. 
Vision-oriented surveys \cite{rahman2024mamba,liu2024vision} further examine the adaptation of Mamba to visual data, including different scanning strategies and applications in medical imaging, remote sensing, and visual recognition. These surveys provide broad coverage of SSM/Mamba architectures and applications but do not specifically focus on their adaptation to graph-structured data.

More recently, Xu et al.~\cite{xu2026graph} presented a dedicated survey of
Graph Mamba, primarily emphasizing its theoretical foundations. Their work
establishes connections between graph propagation and state-space evolution
and interprets Graph Mamba as a selective dynamical framework across different
graph domains. In contrast, our survey adopts an architectural and comparative
perspective. Rather than primarily focusing on the theoretical interpretation
of Graph Mamba, we systematically examine how selective SSMs are adapted to
graphs through graph-to-sequence construction, node ordering and serialization,
graph-aware state-space integration, scanning strategies, and application-specific
architectural designs. We further compare representative models in terms of
their graph settings, design choices, benchmark performance, and computational
characteristics, and analyze practical challenges related to scalability,
order sensitivity, structural preservation, interpretability, and
reproducibility.

Taken together, existing surveys provide strong coverage of GNNs, Graph
Transformers, general SSMs, and the theoretical foundations of Graph Mamba,
but a systematic architectural synthesis remains limited. In particular,
there is still a need to clarify where graph structure enters the selective
state-space pipeline, how serialization and permutation choices affect model
behavior, how reported efficiency should be interpreted beyond the SSM scan
itself, and under which evaluation conditions the claimed advantages of Graph
Mamba are supported. These gaps motivate the scope of the present survey.

\subsection{Research Questions}

To provide a coherent synthesis of the rapidly evolving Graph Mamba
literature, this survey is organized around four research questions:

\begin{itemize}

    \item \textbf{RQ1: How is graph structure incorporated into selective
    state-space modeling?}
    We examine how graph topology influences token construction, node ordering,
    serialization, state-space operators, structural conditioning, and
    graph--SSM feature integration.
    % [REV-R03] Added (red): RQ1 is tied to the four design questions.
    We describe every architecture by its answers to four design questions: tokenization, ordering, coupling, and integration.

    \item \textbf{RQ2: What architectural, structural, and computational
    implications follow from these design choices?}
    We analyze permutation properties, order sensitivity, long-range
    information propagation, scalability, and the computational overhead
    introduced by graph construction and serialization.

    \item \textbf{RQ3: Under which graph settings, tasks, datasets, and
    evaluation conditions are the reported advantages of Graph Mamba
    supported?}
    We examine benchmark practices, dataset characteristics, reported
    performance, and the comparability of results across studies.

    \item \textbf{RQ4: Which unresolved challenges currently limit reliable
    progress in Graph Mamba?}
    We investigate open problems related to robustness, expressive power,
    implementation efficiency, interpretability, pre-training, and future
    state-space backbones.

\end{itemize}


\subsection{Contributions of the Proposed Survey}

The main contributions of this survey are summarized as follows:

\begin{itemize}

    \item We provide a unified architectural perspective on Graph Mamba,
    clarifying how graph information is transformed, serialized, and integrated
    into selective state-space computation.

    \item We introduce a taxonomy based on four design questions:
    tokenization, ordering, coupling, and integration, with input-, dynamics-,
    and operator-level coupling.

    \item We synthesize Graph Mamba applications, structural and computational
    trade-offs, evaluation practices, and reported performance across diverse
    graph-learning settings.

    \item We identify current limitations and formulate open research directions
    related to robustness, scalability, interpretability, reproducibility, and
    future Graph Mamba designs.

\end{itemize}

\subsection{Survey Methodology}
\label{subsec:methodology}

To construct the literature corpus, we searched major academic databases,
including IEEE Xplore, ACM Digital Library, Scopus, Web of Science, and Google
Scholar, using combinations of keywords such as ``Graph Mamba,'' ``Graph State
Space Model,'' ``Mamba graph learning,'' ``selective state space graph,'' and
``state space model graph.'' The search covered studies published from 2023 to
September 2026 and was complemented by backward and forward citation searches
of relevant publications.

We included peer-reviewed articles and relevant preprints that explicitly
integrate Mamba, selective state-space models, or closely related structured
SSMs with graph-structured representations. Studies applying Mamba exclusively
to sequential, image, or other non-graph data without an explicit graph
modeling component were excluded. Earlier structured-SSM approaches were
retained when they provide direct architectural foundations for subsequent
Graph Mamba developments.

The selected studies were analyzed according to their graph setting,
graph--SSM coupling mechanism, serialization or scanning strategy, learning
objective, application domain, benchmark usage, and reported computational and
predictive performance. This organization supports both the architectural
taxonomy developed in this survey and the subsequent cross-domain and
experimental comparisons.
% [REV-R05] Added (red): link between the coding of studies and the four design questions.
%Each architecture was then described by its answers to the four design questions introduced in Section~\ref{sec:unified-framework}.

% [REV-R06] Leftover note "Check at the end!" removed.
% [REV-R06-ORIGINAL] \subsection{Paper Organization} \textcolor{blue}{Check at the end!}
\subsection{Paper Organization}

% [REV-R06/R23] Replaced: the original paragraph referred to undefined labels (sec2, sec3, sec4, sec5, Section6new -> "Section ??"), omitted Section 4, and mentioned "learning strategies", which no section covers.
% [REV-R06-ORIGINAL] The remainder of this paper is organized as follows. Section~\ref{sec2}
% [REV-R06-ORIGINAL] introduces the foundations of GNNs, SSMs, and Mamba. Section~\ref{sec3}
% [REV-R06-ORIGINAL] presents Graph Mamba architectures and learning strategies. Section~\ref{sec4}
% [REV-R06-ORIGINAL] reviews recent applications across different graph settings and domains.
% [REV-R06-ORIGINAL] Section~\ref{sec5} summarizes evaluation practices and benchmarking datasets,
% [REV-R06-ORIGINAL] while Section~\ref{Section6new} provides comparative analyses of Graph Mamba
% [REV-R06-ORIGINAL] and recent alternatives. Section~\ref{sec6} discusses challenges, open issues,
% [REV-R06-ORIGINAL] and future research directions. Finally, Section~\ref{sec8} concludes the paper.
 


The remainder of this paper is organized as follows. Section~\ref{sec:foundations} introduces the foundations of GNNs, SSMs, and Mamba, together with the unified graph-to-state-space formulation and the four design questions. Section~\ref{sec:taxonomy} presents the architectural taxonomy of general-purpose, dynamic and spatio-temporal, and heterogeneous and higher-order Graph Mamba models. Section~\ref{sec:structural-properties} analyzes the structural and computational implications of these design choices. Section~\ref{sec:applications} reviews domain applications and cross-domain design patterns. Section~\ref{sec:evaluation} examines benchmark evidence and reported performance. Section~\ref{sec6} discusses open challenges and future research directions, and Section~\ref{sec8} concludes the paper.
%\begin{figure}[ht!]
%    \centering
    %\includegraphics[width=1.1\linewidth]{Figures/Survey_strc.png}
%    \caption{The structure of the survey }
%    \label{fig:enter-label}
%\end{figure}

% ==========================================================================
% Section 3 -- Foundations and Unified Formulation (replaces Preliminaries and the
% conceptual part of the old "Graph Mamba Architectures" section)
% Notation is unified: u = input sequence, s = latent state, y = output; H = node
% representations, U = token sequence. (The old text used x for both the latent state and
% the input and h for both node representations and the latent state.)
% Labels defined here: sec:foundations, sec:graph-notation, sec:gnn-background,
%   sec:ssm-background, sec:unified-framework, fig:graph-mamba (the architecture figure),
%   eq:equivariance, eq:invariance, eq:mp-layer, eq:zoh-A, eq:zoh-B, eq:recurrence,
%   eq:alpha-retention, eq:ssd, eq:construct ... eq:readout, eq:multi-scan.
% Labels used from other sections: fig:graph-types, fig:serialization, sec:taxonomy,
%   sec:properties, sec:evaluation.
% Needs: amsmath, array, booktabs, colortbl, \input{gm_palette}
% ==========================================================================
\section{Foundations and Unified Formulation}
\label{sec:foundations}

This section introduces the notation, background, and unified framework used to analyze Graph Mamba architectures and graph–SSM interactions.

%---------------------------------------------------------------------------
\subsection{Graph Learning Foundations}
\label{sec:graph-foundations}

Let $G=(V,E)$ denote a graph with $|V|=n$ nodes and $|E|=m$ edges,
node-feature matrix $X\in\mathbb{R}^{n\times D}$, and adjacency matrix
$A\in\mathbb{R}^{n\times n}$. Node representations at layer $l$ are denoted
by $H^{(l)}=[h_1^{(l)},\ldots,h_n^{(l)}]^\top$. Depending on the application,
the graph may be given directly, as in citation or interaction
networks, or constructed from non-graph data such as images, biosignals,
clinical variables, or time series. Graph learning tasks are commonly defined
at the node level, such as node classification or regression; at the edge
level, such as link prediction; or at the graph level, such as graph
classification, regression, or anomaly detection.

GNNs learn graph representations by propagating and
aggregating information over the graph structure. A generic message-passing
layer can be written as

\begin{equation}
h_i^{(l+1)}
=
\sigma\left(
\sum_{j\in\mathcal{N}(i)}
a_{ij}\,
W^{(l)}h_j^{(l)}
\right),
\label{eq:message-passing}
\end{equation}

% [REV-R07] Added (red): Eq. (1) has no self term; clarified that N(i) includes i when self-loops are added.
where $\mathcal{N}(i)$ denotes the neighborhood of node $i$ (including $i$ itself when self-loops are added),
$W^{(l)}$ is a learnable transformation, $a_{ij}$ controls the contribution
of node $j$, and $\sigma(\cdot)$ is a nonlinear activation function.
Different GNN architectures mainly differ in how the neighborhood and
aggregation coefficients are defined. Graph Convolutional Networks (GCNs)
use degree-normalized aggregation \cite{kipf2016semi}, Graph Attention
Networks (GATs) learn attention coefficients between neighboring nodes
\cite{velivckovic2017graph}, and GraphSAGE samples and aggregates local
neighborhoods to support inductive learning on large graphs
\cite{hamilton2017inductive}.

Because message passing propagates information locally, a single layer
typically communicates only across one-hop neighborhoods. Modeling
dependencies between distant nodes therefore requires multiple propagation
layers, which may increase computation and can introduce problems such as
over-smoothing or over-squashing. Graph Transformers provide an alternative
by allowing global interactions through self-attention. However, full
attention over $n$ graph nodes has $\mathcal{O}(n^2)$ time and memory
complexity. Architectures such as GraphGPS
\cite{rampavsek2022recipe} and Exphormer
\cite{shirzad2023exphormer} combine local graph processing with more global
interaction mechanisms. This local--global design provides an important
architectural precedent for Graph Mamba, where selective state-space models
can replace or complement attention for long-range information propagation.

A second fundamental property of graph learning is that graph representations
should not depend on the arbitrary indices assigned to nodes. Let
$P_\sigma$ denote the permutation matrix associated with a node relabeling
$\sigma\in S_n$. A node-level function $f$ is
\emph{permutation-equivariant} if

\begin{equation}
f(P_\sigma X,\,
  P_\sigma A P_\sigma^{\top})
=
P_\sigma f(X,A),
\label{eq:equivariance}
\end{equation}

meaning that relabeling the input nodes produces the corresponding relabeling
of the output node representations. A graph-level function $g$ is
\emph{permutation-invariant} if

\begin{equation}
g(P_\sigma X,\,
  P_\sigma A P_\sigma^{\top})
=
g(X,A),
\label{eq:invariance}
\end{equation}

so that the graph-level prediction remains unchanged under node relabeling.
Conventional message-passing GNNs satisfy these properties when their
% [REV-R07] Terminology: "permutation independent" -> "permutation invariant".
% [REV-R07-ORIGINAL] aggregation operators are permutation independent.
aggregation operators are permutation invariant.

These properties become especially important when state-space models are
adapted to graphs. Unlike graphs, which do not generally possess a unique
canonical ordering, sequential SSMs process representations according to an
explicit sequence. Consequently, Graph Mamba architectures must determine how
graph elements are converted into ordered representations while preserving structural information.  

%---------------------------------------------------------------------------
%---------------------------------------------------------------------------
\subsection{Graph Settings Considered in This Survey}
\label{sec:graph-settings}

Graph Mamba models are applied to graph structures with different structural and
temporal characteristics. Figure~\ref{fig:graph-types} summarizes the main settings considered in
this survey. A static graph has a fixed graph structure during the learning instance,
whereas a dynamic graph evolves over time through changes in nodes, edges, or
their attributes. A homogeneous graph contains a single node and relation type,
while a heterogeneous graph includes multiple node or relation types and
therefore requires type-aware representation learning.
Formally, a heterogeneous graph can be represented as
$G=(V,E,\phi,\psi)$, where $\phi:V\rightarrow\mathcal{A}$ maps nodes to
node types and $\psi:E\rightarrow\mathcal{R}$ maps edges to relation types.
Here, type information is distinct from ordinary node or edge features because
different types may require different representation spaces or type-specific
parameters. A spatio-temporal graph combines graph-based spatial relationships with
observations that vary over time. Finally, higher-order graphs extend pairwise
edges to interactions involving multiple entities, such as simplices or cells
of different ranks. These characteristics are not mutually exclusive. For example, a real-world
graph may be heterogeneous, dynamic, and spatio-temporal at the same time.
Accordingly, they are used here as broad graph settings rather than as strict
exclusive categories.

\begin{figure*}[h]
    \centering
    \resizebox{0.9\textwidth}{!}{%
        \input{Figures/fig_graph_types_tikz}
    }
    % [REV-NOTE] The caption mentions higher-order structures; check that Figures/fig_graph_types_tikz contains a higher-order panel (the compiled v9 figure shows five panels without it).
    \caption{Graph settings considered in this survey. Static graphs have fixed structure; dynamic graphs gain or lose nodes and edges over time; homogeneous graphs have one node and relation type, whereas heterogeneous graphs have several; spatio-temporal graphs attach time-varying signals to spatially anchored nodes. The settings are not mutually exclusive, and real-world graphs frequently combine them.}
    \label{fig:graph-types}
\end{figure*}

%---------------------------------------------------------------------------
\subsection{State-Space Models and Mamba}
\label{sec:ssm-background}

State-Space Models (SSMs) provide a mathematical framework for modeling
sequential data through a latent state that evolves over time. In continuous
form, an SSM maps an input signal $u(t)$ to an output $y(t)$ through



\begin{equation}
s'(t)=A\,s(t)+B\,u(t),\qquad y(t)=C\,s(t)+D\,u(t),
\label{eq:ssm-continuous}
\end{equation}

where $s(t)$ denotes the latent state, $A$ governs the state dynamics,
$B$ maps the input into the state, $C$ projects the state to the output,
and $D$ represents a direct input-to-output connection. For discrete
sequences, the continuous system is discretized using a step size $\Delta$.
Under zero-order-hold discretization, the corresponding parameters are

\begin{equation}
\bar A=\exp(\Delta A),\qquad \bar B=(\Delta A)^{-1}\big(\exp(\Delta A)-I\big)\Delta B,
\label{eq:zoh-A}
\end{equation}


resulting in the recurrence

\begin{equation}
s_t=\bar{A}s_{t-1}+\bar{B}u_t,
\qquad
y_t=C s_t+D u_t.
\label{eq:ssm-recurrence}
\end{equation}

Because the step size $\Delta$ is explicit, the same continuous-time system can
be discretized at non-uniform intervals. Continuous-time dynamic-graph models
can exploit this property by deriving $\Delta$ from the elapsed time between
successive events.

Modern SSM architectures progressively modify this formulation to improve
long-range modeling and computational efficiency.

\paragraph{\textbf{Structured State-Space Models.}}
S4 introduces a structured parameterization of the state matrix that enables
efficient modeling of long sequences \cite{gu2021efficiently}. Its design is
closely related to the High-order Polynomial Projection Operator (HiPPO),
which provides a principled mechanism for compressing the history of a
continuous input signal into a finite-dimensional state
\cite{gu2020hippo}. After discretization, the recurrence can be represented as
a convolution with a structured kernel. S5 further simplifies this
framework through a multi-input, multi-output formulation that supports
efficient parallel sequence processing \cite{smith2022simplified}. These
structured SSMs provide an important precursor to Graph Mamba because they
demonstrate how long-range dependencies can be modeled without quadratic
self-attention.

\paragraph{\textbf{Selective State-Space Modeling with Mamba.}}
A limitation of earlier structured SSMs is that their state-space parameters
are largely independent of the current input. Mamba addresses this limitation
through the Selective State Space Model (S6), in which key parameters become
input dependent \cite{gu2023mamba}. For an input token $u_t$, the selection
mechanism can be written as

\begin{equation}
B_t=f_B(u_t),
\qquad
C_t=f_C(u_t),
\qquad
\Delta_t=
\mathrm{softplus}\!\left(
b_{\Delta}+f_{\Delta}(u_t)
\right),
\label{eq:mamba-selection}
\end{equation}

where $f_B$, $f_C$, and $f_{\Delta}$ are learnable input-dependent
projections and $b_{\Delta}$ is a learnable bias.
% [REV-R08] Added (red): Eq. (8) uses bar A_t and bar B_t, which were never defined. [REV-NOTE] verify wording against Gu and Dao (2023), Sec. 3.
Applying the discretization of Equation~\eqref{eq:zoh-A} with the input-dependent step size gives $\bar{A}_t=\exp(\Delta_t A)$, while Mamba simplifies the input matrix to $\bar{B}_t\approx\Delta_t B_t$. The discretized parameters
therefore vary across positions, leading to the selective recurrence

\begin{equation}
s_t
=
\bar{A}_t s_{t-1}
+
\bar{B}_t u_t,
\qquad
y_t
=
C_t s_t
+
D u_t .
\label{eq:mamba-recurrence}
\end{equation}

Unlike fixed SSMs, Mamba can dynamically regulate how information is written
to, retained in, and read from the latent state. The recurrence is evaluated
through a hardware-aware selective scan whose computational cost scales
linearly with sequence length.

The contribution of an earlier input token $u_k$ to a later output $y_t$ can
be exposed by unrolling the recurrence:
\begin{equation}
y_t=\sum_{k=1}^{t}\alpha_{t\leftarrow k}\,u_k+D\,u_t,\qquad
\alpha_{t\leftarrow k}=C_t\Big(\prod_{j=k+1}^{t}\bar A_j\Big)\bar B_k ,
\label{eq:alpha-retention}
\end{equation}

The coefficient $\alpha_{t\leftarrow k}$ can be interpreted as an
input-dependent retention profile describing how strongly information from
position $k$ contributes to the output at position $t$. 
In contrast to node representations in message-passing GNNs, the SSM state $s_t$
summarizes the tokens processed up to position $t$ in a fixed-size memory.
Moreover, the input-dependent step size $\Delta_t$ controls how much previous
state information is retained and how strongly the current token is written.


\paragraph{\textbf{Mamba-2 and State Space Duality.}}
Mamba-2 extends selective state-space modeling through State Space Duality
(SSD), which establishes a close connection between structured SSMs and
attention-like matrix formulations \cite{dao2024transformers}. Mamba-2
simplifies the transition structure and enables more efficient use of modern
matrix-multiplication hardware while supporting larger state dimensions. This
results in substantially faster computation in many sequence-modeling settings.

Although most existing Graph Mamba architectures are based on the original
Mamba selective scan, Mamba-2 and SSD are relevant to future graph models
because they provide an alternative sequence backbone and expose structured
sequence interactions more explicitly.

Importantly, none of these SSM formulations directly encode graph topology.
They operate on an ordered sequence of representations. Therefore, adapting
Mamba to graph-structured data requires an additional mechanism that determines
what constitutes a token, how graph elements are ordered or scanned, how graph
structure conditions the state-space computation, and how the resulting
sequence outputs are mapped back to the graph domain. 

%---------------------------------------------------------------------------

%---------------------------------------------------------------------------
\subsection{Unified Graph-to-State-Space Formulation}
\label{sec:unified-framework}

The Graph Mamba architectures reviewed in this survey can be interpreted
through a common processing pipeline that connects graph construction,
graph-derived representations, sequence formation, selective state-space
processing, and graph-level output generation. This formulation makes explicit
where graph structure enters the computation and provides a common reference
for comparing different architectural designs.
Figure~\ref{fig:pipeline} illustrates this pipeline and shows where the four design questions—tokenization, ordering, coupling, and integration—arise.

% ---------- FIGURE 1 HERE ----------
\begin{figure*}[h]
\centering

\resizebox{0.99\textwidth}{!}{%
    \input{Figures/fig_pipeline_tikz}}
    \caption{Unified graph-to-state-space pipeline. A graph is constructed or given (1), converted into tokens (2, D1), ordered and scanned (3, D2), processed by a selective SSM whose parameters may be conditioned on graph structure (4, D3), and mapped back to graph elements and fused with an optional local message-passing branch (5, D4). Labels E1–E5 mark where graph structure enters each stage. }
\label{fig:pipeline}
\end{figure*}

% -----------------------------------


Let $\mathcal{D}$ denote the raw input data and let
$G=(V,E,X)=\mathcal{C}(\mathcal{D}),$
where $\mathcal{C}(\cdot)$ denotes the graph-construction procedure. For
native graph datasets, this operation is the identity. In application-specific
settings, however, the graph may be derived from images, biosignals, clinical
variables, time series, or other structured observations.

An optional graph encoder first produces local graph-aware representations,

\begin{equation}
H_{\mathrm{local}}
=
\Phi_{\mathrm{graph}}(X,A),
\label{eq:local-graph-representation}
\end{equation}

where $\Phi_{\mathrm{graph}}$ may correspond to a GCN, GAT, GraphSAGE, or
another graph-processing operator. Architectures without a separate local
branch may instead construct tokens directly from the input graph.

A tokenization function $\mathcal{T}$ then transforms graph elements or
graph-derived representations into a set of state-space inputs,

\begin{equation}
U
=
\mathcal{T}(G,H_{\mathrm{local}})
=
[u_1,\ldots,u_T]^\top
\in\mathbb{R}^{T\times D}.
\label{eq:graph-tokenization}
\end{equation}

Depending on the architecture, a token may represent an individual node, a
sampled neighborhood, a graph walk, a metapath instance, a higher-order cell,
a graph snapshot, or another structured graph component.

Because selective SSMs operate on ordered inputs, the graph-derived tokens
must either be arranged according to an ordering $\pi$ or processed by an
operator that avoids explicit serialization. For serialization-based models,

\begin{equation}
U_{\pi}
=
P_{\pi}U,
\label{eq:serialization}
\end{equation}

where $P_{\pi}$ is the permutation matrix associated with the selected
ordering. The ordering may be determined from structural importance,
neighborhood traversal, temporal information, node or relation type,
higher-order rank, or other graph-dependent criteria.

The resulting sequence is processed by a selective state-space model,

\begin{equation}
S_t
=
\bar{A}_t S_{t-1}
+
\bar{B}_t U_{\pi,t},
\qquad
Y_t
=
C_t S_t
+
D U_{\pi,t},
\label{eq:graph-mamba-state}
\end{equation}

where the state-space parameters may depend only on the current token, as in
standard Mamba, or may additionally be conditioned on graph-derived structural
information. 

In the latter case, for structural context $z_t$,

\begin{equation}
B_t=f_B(U_{\pi,t},z_t),\qquad
C_t=f_C(U_{\pi,t},z_t),\qquad
\Delta_t=
\mathrm{softplus}
\left(
b_{\Delta}+f_{\Delta}(U_{\pi,t},z_t)
\right).
\label{eq:graph-conditioned-selection}
\end{equation}

After state-space processing, sequence outputs are mapped back to their
corresponding graph elements. For node-level serialization, this can be
expressed as

\begin{equation}
H_{\mathrm{global}}
=
P_{\pi}^{\top}Y,
\label{eq:graph-restoration}
\end{equation}

whereas more general tokenizations require an aggregation or scattering
operator that maps token-level outputs back to nodes, edges, or higher-order
graph elements.

Finally, architectures that combine local graph processing and global
state-space modeling fuse both representations,

\begin{equation}
H
=
\Phi_{\mathrm{fuse}}
\left(
H_{\mathrm{local}},
H_{\mathrm{global}}
\right),
\qquad
\hat{y}
=
\rho(H),
\label{eq:graph-mamba-fusion}
\end{equation}

where $\Phi_{\mathrm{fuse}}$ may implement addition, concatenation, gating, or
a learned projection, and $\rho(\cdot)$ denotes the task-specific prediction
head.

% [REV-R09] Replaced: the five "principal points" competed with the D1--D4 taxonomy of Section 3 and with the list in the conclusion.
% The formulation is now summarized by four design questions (+ graph construction as context), used verbatim in every section.
% [REV-R09-ORIGINAL] This unified formulation identifies five principal points at which graph
% [REV-R09-ORIGINAL] structure can influence a Graph Mamba architecture:
% [REV-R09-ORIGINAL] 
% [REV-R09-ORIGINAL] \begin{enumerate}
% [REV-R09-ORIGINAL]     \item \textbf{Graph construction:} which entities, relations, and features
% [REV-R09-ORIGINAL]     define the graph;
% [REV-R09-ORIGINAL] 
% [REV-R09-ORIGINAL]     \item \textbf{Token construction:} what graph element or structural context
% [REV-R09-ORIGINAL]     each state-space input represents;
% [REV-R09-ORIGINAL] 
% [REV-R09-ORIGINAL]     \item \textbf{Serialization and scanning:} how graph-derived tokens are
% [REV-R09-ORIGINAL]     ordered and traversed;
% [REV-R09-ORIGINAL] 
% [REV-R09-ORIGINAL]     \item \textbf{State-space conditioning:} whether topology, time, type,
% [REV-R09-ORIGINAL]     rank, or other graph information directly modifies the selective
% [REV-R09-ORIGINAL]     state-space computation;
% [REV-R09-ORIGINAL] 
% [REV-R09-ORIGINAL]     \item \textbf{Fusion and readout:} how state-space outputs are aligned with
% [REV-R09-ORIGINAL]     the graph and combined with local graph representations.
% [REV-R09-ORIGINAL] \end{enumerate}
% [REV-R09-ORIGINAL] 
% [REV-R09-ORIGINAL] Existing Graph Mamba models differ primarily in which of these stages are
% [REV-R09-ORIGINAL] graph-aware and in how strongly graph structure is coupled with the selective
% [REV-R09-ORIGINAL] SSM. These differences form the basis of the architectural taxonomy developed
% [REV-R09-ORIGINAL] in the following section.
The pipeline raises four design questions, which form the taxonomy used
throughout this survey. \emph{Tokenization} concerns what each state-space
input represents (Equation~\eqref{eq:graph-tokenization}). \emph{Ordering}
concerns how tokens are arranged and scanned
(Equation~\eqref{eq:serialization}). \emph{Coupling} concerns how deeply
graph structure enters the state dynamics. At the \textbf{input level}, the
graph shapes only the tokens and their order. At the \textbf{dynamics level},
graph information conditions the selective parameters, the discretization, or
the transition (Equation~\eqref{eq:graph-conditioned-selection}). At the
\textbf{operator level}, the state-space operator itself is defined on the
graph. \emph{Integration} concerns how the state-space branch is combined with
local graph propagation (Equations~\eqref{eq:graph-restoration}
and~\eqref{eq:graph-mamba-fusion}): the model may run standalone, follow or
precede a graph encoder sequentially, run in parallel with it, or be embedded
inside graph propagation.
Graph construction is not a design
question of the architecture but part of its context. When the graph is
derived from other data, its construction becomes part of the method and must
be considered when results are compared.



\section{Taxonomy of Graph Mamba Architectures}
\label{sec:taxonomy}
This section addresses RQ1 by examining how graph structure is incorporated into selective state-space modeling. We organize existing Graph Mamba architectures according to four complementary design dimensions—tokenization, ordering, coupling, and integration—and analyze how these choices vary across static, dynamic, spatio-temporal, heterogeneous, and higher-order graph settings.

\subsection{Taxonomy Framework}
\label{subsec:taxonomy-dimensions}

Building on the unified formulation in
Section~\ref{sec:unified-framework}, we classify Graph Mamba architectures
according to four complementary dimensions that capture where graph structure
interacts with selective state-space processing. 
Table~\ref{tab:taxonomy-dimensions} summarizes the four dimensions.%, while Figure~\ref{fig:serialization} illustrates representative token construction, serialization, and scanning strategies associated mainly with D1 and D2.



\begin{table*}[h]
\centering
% [REV-R10] Table 1 aligned with the four design questions (names added in red; original text kept).
\caption{
Taxonomy dimensions used to classify Graph Mamba architectures.
}
\label{tab:taxonomy-dimensions}
\small
\renewcommand{\arraystretch}{1.25}
\resizebox{\textwidth}{!}{
\begin{tabular}{@{}p{1cm}p{3cm}p{3cm}p{7.0cm}@{}}
\toprule
\textbf{Dim.} &
\textbf{Architectural Aspect} &
\textbf{Main Question} &
\textbf{Representative Design Choices} \\
\midrule

D1 &
\textbf{Tokenization:} Graph representation and token construction &
What information is presented to the state-space model? &
Nodes, neighborhoods, walks, metapaths, spectral components,
snapshots, higher-order cells \\

D2 &
\textbf{Ordering:} Serialization, scanning, and directionality &
How are graph-derived representations ordered and traversed? &
Structural ranking, traversal, temporal ordering, type-aware or
rank-aware ordering, uni-/bi-/multi-directional scanning \\

D3 &
\textbf{Coupling:} Structural conditioning of state-space computation &
Does graph information directly affect the selective state dynamics? &
Graph-conditioned $B_t$, $C_t$, or $\Delta_t$; modified transition;
graph-aware state-space operator. Levels: input-level, dynamics-level, operator-level \\

D4 &
\textbf{Integration:} Integration with graph propagation or attention &
How is the SSM combined with other graph-learning mechanisms? &
Parallel local--global branches, graph encoder followed by Mamba,
embedded state evolution, attention--SSM hybrids. Types: standalone, sequential, parallel, embedded. Attention--SSM hybrids are classified according to how the SSM branch is integrated within the architecture. \\

\bottomrule
\end{tabular}}
\end{table*}

\begin{figure*}[h]
\centering
\resizebox{0.99\textwidth}{!}{%
    \input{Figures/fig_graph_serialization_updated}
}
\caption{
Representative token construction, serialization, and scanning strategies in
Graph Mamba architectures. Panels (a)--(e) illustrate representative D1 and
D2 mechanisms across structural, neighborhood-based, heterogeneous, temporal,
and higher-order graph settings. The bottom panel compares unidirectional,
bidirectional, and multi-order scanning.}
\label{fig:serialization}
\end{figure*}

The four dimensions are not mutually exclusive. For example, a model may use neighborhood-based tokenization (D1), bidirectional structure-aware scanning (D2), and a parallel message-passing branch (D4), while leaving the Mamba selection mechanism unchanged and therefore coupling graph information only at the input level (D3). This multi-label view enables architectures developed for different graph settings to be compared using the same criteria.  % [REV-R12] Names of the design questions added next to D1/D2 (red).
As illustrated in Figure~\ref{fig:serialization}, D1, tokenization, determines the graph-derived
units presented to the SSM, whereas D2, ordering, determines how these units are ordered
and scanned.

For presentation, we organize the reviewed architectures into three broad
graph settings: static and general-purpose graphs, dynamic and spatio-temporal
graphs, and heterogeneous and higher-order graphs. Within each group, models are compared according to the four
architectural dimensions D1--D4.
% [REV-R11] Added (red): reading guide. Each model is described by its answers to the four questions in the same order, and each graph setting emphasizes a different question (this finding is repeated in the conclusion).
Each architecture below is described by its answers to the four design
questions, in the same order: what it reads (tokenization, D1), how it orders
what it reads (ordering, D2), at which level it couples graph structure to the
state dynamics (coupling, D3), and how it integrates the state-space model
with graph propagation (integration, D4). The three graph settings emphasize
different questions. In static graphs the decisive question is ordering,
because nodes have no natural sequence. In dynamic graphs time largely answers
the ordering question, so design effort shifts to coupling. In heterogeneous
and higher-order graphs, most structure is carried by tokenization.


\subsection{Static and General-Purpose Graph Models}
\label{subsec:static-models}

Static graphs provide a particularly challenging setting for Graph Mamba
because, unlike temporal data, their nodes do not possess an intrinsic
sequential order. Existing approaches therefore differ mainly in how they
adapt graph structure to state-space processing. %Representative models include GSSC~\cite{huang2024can}, GMN~\cite{behrouz2024graph}, Graph-Mamba~\cite{wang2024graph}, GrassNet~\cite{zhao2024grassnet}, MbaGCN~\cite{he2025mamba}, DMbaGCN~\cite{he2026dual}, and GLADMamba~\cite{fu2025gladmamba}. 

\paragraph{\textbf{Tokenization and serialization.}}
Graph-Mamba~\cite{wang2024graph} and GMN~\cite{behrouz2024graph}
illustrate two distinct strategies for adapting unordered graphs to sequential
state-space processing. Graph-Mamba operates on node representations and
introduces graph-centric node prioritization before applying a Graph-Mamba
Block within the GraphGPS framework. In its default configuration, nodes are
prioritized by degree, while nodes with the same priority are randomly
permuted during training to reduce ordering bias. The Graph-Mamba Block
replaces the global attention component of GraphGPS, while a Message Passing Neural Network (MPNN) branch
provides local graph processing. Its main graph-aware mechanisms therefore
correspond to serialization and scanning (D2) and local--global graph--SSM
integration (D4).
% [REV-R13] Added (red): closing sentence answering the four design questions in order. [REV-NOTE] verify the coupling level / integration type against the original paper.
%In terms of the four design questions, Graph-Mamba reads node representations, orders them by degree with random permutation among ties, couples the graph at the input level, and runs in parallel with a message-passing branch inside GraphGPS.
GMN~\cite{behrouz2024graph}, in contrast, places greater emphasis on graph
tokenization. Its framework explicitly comprises neighborhood tokenization,
token ordering, local encoding, and a bidirectional selective SSM encoder.
Sampled neighborhood structures are transformed into graph tokens, which are
subsequently processed in both scanning directions. Thus, graph structure is
introduced primarily through token construction (D1) and ordering and
bidirectional scanning (D2). GMN can additionally employ an optional
message-passing module to augment the state-space representation with local
graph inductive bias (D4).
% [REV-R13] Added (red): closing sentence answering the four design questions in order.
%In terms of the four design questions, GMN reads sampled neighborhood tokens, orders them and scans them in both directions, couples the graph at the input level, and optionally adds a parallel message-passing branch.

\paragraph{\textbf{Order-free and spectral state-space processing.}}
GSSC~\cite{huang2024can} avoids explicit graph-to-sequence serialization.
It extends state-space convolution to graphs using global
permutation-equivariant set aggregation together with factorizable graph
kernels based on structural and relative positional information. As a result,
the model provides global information propagation while remaining
permutation equivariant. GSSC therefore represents an operator-level strategy
in which graph structure is incorporated directly into the state-space
convolution rather than through a selected node ordering.
% [REV-R13] Added (red): closing sentence answering the four design questions in order.
%In terms of the four design questions, GSSC needs no ordering, couples the graph at the operator level because the state-space convolution is defined on the graph, and runs standalone.
GrassNet~\cite{zhao2024grassnet} follows a spectral strategy. It performs
spectral decomposition of the graph and constructs spectrum embeddings that
are ordered according to graph frequencies. Its SSM graph filter then scans
the resulting spectral sequence in both directions, allowing each frequency
component to interact with higher- and lower-frequency components. GrassNet
therefore combines spectral representation (D1), frequency-based
bidirectional scanning (D2), and SSM-based spectral filtering (D3).
% [REV-R13] Added (red): closing sentence answering the four design questions in order.
%In terms of the four design questions, GrassNet reads spectral components, orders them by graph frequency and scans them in both directions, couples the graph at the operator level because the state-space model acts as a spectral graph filter, and runs standalone.

\paragraph{\textbf{State-space dynamics within graph propagation.}}
MbaGCN~\cite{he2025mamba} integrates the Mamba paradigm directly into graph
convolution rather than first serializing the graph into a node sequence. Its
architecture contains three main components: a Message Aggregation Layer
(MAL), a Selective State Space Transition Layer (S3TL), and a Node State
Prediction Layer (NSPL). MAL aggregates neighborhood information, whereas
S3TL selectively combines neighborhood features with node-specific
representations and regulates information obtained from neighborhoods of
different orders. NSPL further controls information flow within the same-order
neighborhood. Thus, MbaGCN primarily couples selective state evolution (D3)
with graph message propagation (D4).
% [REV-R13] Added (red): closing sentence answering the four design questions in order.
%In terms of the four design questions, MbaGCN reads neighborhood representations of increasing order, orders them by hop order, couples the graph at the dynamics level through a selective transition that regulates aggregation, and is embedded inside graph propagation.
DMbaGCN~\cite{he2026dual} extends this idea by addressing graph
representation learning from complementary local and global perspectives. Its
Local State-Evolution Mamba (LSEMba) first aggregates local neighborhood
information through stacked graph convolution layers and then models the
progressive evolution of node representations across layers using selective
state-space modeling. Its Global Context-Aware Mamba (GCAMba) processes a
sequence of node representations using bidirectional Mamba to capture
graph-wide dependencies. The two representations are subsequently combined,
providing both local state evolution and global contextual modeling.
% [REV-R13] Added (red): closing sentence answering the four design questions in order. [REV-NOTE] verify the coupling level / integration type against the original paper.
%In terms of the four design questions, DMbaGCN reads layer-wise node states and a global node sequence, orders them by layer depth and by a bidirectional global scan, couples the graph at the dynamics level, and combines an embedded local branch with a parallel global branch.

\paragraph{\textbf{Multi-view and spectrum-guided integration.}}
GLADMamba~\cite{fu2025gladmamba} adapts selective state-space modeling to
unsupervised graph-level anomaly detection. Its View-Fused Mamba (VFM)
integrates information from multiple graph views, while its Spectrum-Guided
Mamba (SGM) introduces explicit spectral information through the Rayleigh
quotient. In particular, the spectral representation is used to parameterize
the selective SSM parameters $B$, $C$, and $\Delta$, thereby directly
conditioning state-space updates on graph spectral information. GLADMamba
therefore combines multi-view graph representation (D1), direct
spectrum-guided state-space conditioning (D3), and multi-view fusion (D4).


In summary, static and general-purpose Graph Mamba designs address the lack of a natural node order through different mechanisms, including structural serialization, neighborhood-based tokenization, spectral representations, and tighter integration of state-space processing with graph propagation. These approaches reveal that the key architectural distinction lies in where graph structure enters the state-space pipeline and how strongly it influences the resulting computation.
 


\subsection{Dynamic and Spatio-Temporal Graph Models}
\label{subsec:dynamic-models}

Dynamic and spatio-temporal graphs introduce an explicit temporal dimension
that distinguishes them from static graph settings. In many cases, time
provides a natural ordering of graph snapshots, interactions, or observations,
so the main challenge is not simply how to serialize graph elements, but how
to jointly represent evolving topology and temporal dependencies.
%Representative approaches include GraphSSM~\cite{li2024state}, GSSM~\cite{zhou2024graph}, DG-Mamba~\cite{yuan2025dg}, DyGMamba~\cite{ding2024dygmamba}, DyG-Mamba~\cite{li2024dyg}, STG-Mamba~\cite{li2024stg}, SpoT-Mamba~\cite{choi2024spot},
% [REV-R14] Missing citations added (red). PS-Mamba~\cite{dong2025ps}, FuzzMamba~\cite{chen2026modeling}, and STMAGRN~\cite{zhang2025spatio}. These architectures mainly differ in whether graph structure influences the temporal representation, the scanning strategy, the state-space dynamics, or a separate graph-processing component.

\paragraph{\textbf{Graph-aware temporal state evolution.}}
GraphSSM~\cite{li2024state} develops a state-space formulation specifically
for discrete-time temporal graphs, which are represented as sequences of graph
snapshots. Its GHIPPO formulation incorporates structural information through
Laplacian regularization, while a mixed discretization mechanism accounts for
graph changes occurring between observed snapshots. GraphSSM therefore
integrates evolving topology directly into state evolution rather than first
reducing the graph to an arbitrary node sequence.
% [REV-R13] Added (red): closing sentence answering the four design questions in order.
%In terms of the four design questions, GraphSSM reads graph snapshots in temporal order, couples the graph at the operator level through Laplacian-regularized state evolution, and runs standalone.
GSSM~\cite{zhou2024graph} provides a related state-space approach for wireless
traffic prediction. It applies graph convolution over fixed and adaptive
base-station graphs and uses the resulting graph representations to estimate
the parameters of an SSM. Since this formulation predates the
input-dependent selection mechanism of Mamba, it is better viewed as a
graph--SSM precursor that demonstrates how spatial graph information can
parameterize temporal state-space modeling.
% [REV-R13] Added (red): closing sentence answering the four design questions in order.
%In terms of the four design questions, GSSM reads base-station time series in temporal order, couples the graph at the dynamics level because graph representations estimate the (non-selective) SSM parameters, and integrates graph convolution sequentially before the SSM.

\paragraph{\textbf{Dynamic interactions and evolving topology.}}
DG-Mamba~\cite{yuan2025dg} addresses dynamic graph structure learning by
combining dynamic message passing with selective state-space modeling across
successive graph states. Structural information from consecutive graph states
is incorporated into the state-space discretization, while an additional
self-supervised information objective regularizes the learned dynamic
structure. Its design therefore combines temporal graph representations (D1),
temporal state-space processing (D2), structure-aware selective dynamics (D3),
and dynamic graph propagation (D4).
% [REV-R13] Added (red): closing sentence answering the four design questions in order. [REV-NOTE] verify the coupling level / integration type against the original paper.
%In terms of the four design questions, DG-Mamba reads successive graph states in temporal order, couples the graph at the dynamics level because structure enters the discretization, and runs in parallel with dynamic message passing.
DyGMamba~\cite{ding2024dygmamba} focuses on continuous-time dynamic graphs and
uses two complementary state-space components. A node-level Mamba SSM encodes
long sequences of historical node interactions, while a time-level Mamba SSM
models edge-specific temporal patterns. The learned temporal patterns are then
used to identify the most relevant information in the interaction history.
Because historical interactions are processed chronologically, temporal order
provides the principal serialization mechanism.
% [REV-R13] Added (red): closing sentence answering the four design questions in order. [REV-NOTE] verify the coupling level / integration type against the original paper.
%In terms of the four design questions, DyGMamba reads historical interactions in chronological order, couples the graph at the input level while a time-level SSM selects relevant history, and runs standalone.
DyG-Mamba~\cite{li2024dyg} also targets continuous-time dynamic graphs, but
addresses irregular temporal intervals more directly. It introduces a
timespan-informed continuous SSM in which the elapsed time between consecutive
events controls the state-space step size and hence the decay of historical
information. The model further uses input-dependent state-space parameters to
selectively retain useful historical information and suppress noisy events.
Thus, whereas DyGMamba separates interaction-history encoding and temporal
pattern modeling, DyG-Mamba explicitly incorporates irregular timespans into
the state transition.
% [REV-R13] Added (red): closing sentence answering the four design questions in order.
%In terms of the four design questions, DyG-Mamba reads timestamped interactions in irregular chronological order, couples the graph at the dynamics level because elapsed time sets the step size $\Delta_t$, and runs standalone.

\paragraph{\textbf{Spatio-temporal modeling and graph-guided scanning.}}
STG-Mamba~\cite{li2024stg} combines a Spatial--Temporal Selective State Space
Module with Kalman Filtering Graph Neural Networks (KFGN). The graph component
models spatial relationships and updates graph-aware representations, while
the selective state-space module captures temporal dependencies at different
granularities. The architecture therefore combines graph-based spatial
processing with selective temporal state evolution, corresponding primarily to
structural conditioning (D3) and graph--SSM integration (D4).
% [REV-R13] Added (red): closing sentence answering the four design questions in order. [REV-NOTE] verify the coupling level / integration type against the original paper.
%In terms of the four design questions, STG-Mamba reads graph-aware node time series in temporal order at multiple granularities, couples the graph at the dynamics level, and runs in parallel with a Kalman-filtering graph network.
SpoT-Mamba~\cite{choi2024spot} introduces a different strategy by constructing
multiple graph-derived spatial sequences using breadth-first search,
depth-first search, and random walks. These complementary traversal sequences
are processed by Mamba to encode spatial neighborhood information, while
temporal state-space processing captures long-range dependencies across time.
% [REV-R15] Dangling list fixed: comma replaced by "with" (red).
% [REV-R15-ORIGINAL] SpoT-Mamba therefore combines graph-derived token construction (D1),
% [REV-R15-ORIGINAL] multi-way spatial scanning and temporal ordering (D2).
SpoT-Mamba therefore combines graph-derived token construction (D1) with
multi-way spatial scanning and temporal ordering (D2).
% [REV-R13] Added (red): closing sentence answering the four design questions in order.
%In terms of the four design questions, SpoT-Mamba reads graph-derived walk sequences, orders them through multiple spatial traversals and a temporal scan, couples the graph at the input level, and runs standalone.
PS-Mamba~\cite{dong2025ps} introduces graph-guided spatial--temporal scanning for human pose
sequence refinement. Its Spatial--Temporal Graph State Space (ST-GSS) block
contains a graph branch with graph convolution, temporal convolution, and a
dynamic graph-weight matrix, together with a Spatial--Temporal State Space
Model. A Graph-guided Spatial--Temporal Scanning strategy constructs four
bidirectional spatial--temporal sequences for state-space processing. The
model therefore integrates graph-guided scanning (D2) with adaptive
joint-interaction modeling and graph--SSM fusion (D4).
% [REV-R13] Added (red): closing sentence answering the four design questions in order.
%In terms of the four design questions, PS-Mamba reads human joints across pose sequences, orders them through four graph-guided bidirectional scans, couples the graph at the input level, and fuses a parallel graph branch within each ST-GSS block.

FuzzMamba~\cite{chen2026modeling} introduces adaptive node-centric sequence construction for
spatio-temporal traffic modeling. It first determines spatial neighborhood
orderings from node distances and then dynamically adjusts these orderings
according to time-varying environmental information. Fuzzy role
representations and environmental disturbance signals are subsequently
combined with neighborhood features to guide Mamba-based sequence modeling.
Its principal contributions therefore lie in adaptive graph-derived
serialization (D2) and the use of dynamic contextual information to guide
selective state-space processing (D3).
% [REV-R13] Added (red): closing sentence answering the four design questions in order. [REV-NOTE] verify the coupling level / integration type against the original paper.
%In terms of the four design questions, FuzzMamba reads node-centric neighborhoods, orders them adaptively by distance as the environment changes, couples the graph at the dynamics level because fuzzy roles and environmental signals guide selection, and runs standalone.
STMAGRN~\cite{zhang2025spatio}  follows a sequential graph--SSM integration strategy for traffic
prediction. A Mamba encoder first extracts long-range spatio-temporal features
from traffic sequences. A Spatio-Temporal Memory (STM) module then identifies
representative traffic patterns and generates dynamic graph embeddings, which
are finally processed by a graph convolutional recurrent decoder combining
GCN and recurrent modeling. Unlike models in which the graph directly defines
the Mamba scan, STMAGRN uses Mamba-based feature extraction before dynamic
graph construction and graph-based decoding, making graph--SSM integration
(D4) its dominant architectural characteristic.
% [REV-R13] Added (red): closing sentence answering the four design questions in order.
%In terms of the four design questions, STMAGRN reads traffic sequences in temporal order, couples the graph at the input level, and integrates sequentially: a Mamba encoder precedes memory-based graph generation and a GCN--RNN decoder.


Taken together, dynamic and spatio-temporal Graph Mamba models emphasize temporal state evolution more strongly than static designs. Because time often provides a natural ordering, the main architectural distinction shifts toward how evolving graph structure influences the state-space computation. In this setting, graph information may remain at the input level, condition the selective dynamics, or be incorporated directly into the state-space operator.

\begin{table*}[h]
\centering
\caption{
Unified taxonomy of representative Graph Mamba architectures according to
tokenization (D1), ordering (D2), coupling (D3), and integration (D4).
}
\label{tab:unified-taxonomy}

\scriptsize
\renewcommand{\arraystretch}{1.15}
\setlength{\tabcolsep}{3pt}

\resizebox{\textwidth}{!}{%
\begin{tabular}{@{}
p{2.2cm}
p{3.6cm}
p{3.6cm}
ccc
cccc
@{}}
\toprule

\multirow{2}{*}{\textbf{Model}} &
\multirow{2}{*}{\textbf{Tokenization (D1)}} &
\multirow{2}{*}{\textbf{Ordering (D2)}} &
\multicolumn{3}{c}{\textbf{Coupling (D3)}} &
\multicolumn{4}{c}{\textbf{Integration (D4)}} \\

\cmidrule(lr){4-6}
\cmidrule(lr){7-10}

& & &
\rotatebox{90}{Input} &
\rotatebox{90}{Dynamics} &
\rotatebox{90}{Operator} &
\rotatebox{90}{Standalone} &
\rotatebox{90}{Sequential} &
\rotatebox{90}{Parallel} &
\rotatebox{90}{Embedded} \\

\midrule

% =========================================================
\multicolumn{10}{l}{\textbf{Static and General-Purpose Graph Models}} \\
\midrule

GSSC~\cite{huang2024can} &
Node features and graph positional information &
Order-free &
& & \cmark &
\cmark & & & \\

GMN~\cite{behrouz2024graph} &
Neighborhood/subgraph tokens &
Ordered, bidirectional &
\cmark & & &
& & \cmark & \\

Graph-Mamba~\cite{wang2024graph} &
Node representations &
Degree-based prioritization and permutation &
\cmark & & &
& & \cmark & \\

GrassNet~\cite{zhao2024grassnet} &
Graph spectrum embeddings &
Frequency-ordered, bidirectional &
& & \cmark &
\cmark & & & \\

MbaGCN~\cite{he2025mamba} &
Neighborhood representations across layers &
No explicit node serialization &
& \cmark & &
& & & \cmark \\

DMbaGCN~\cite{he2026dual} &
Layer-wise node states and node sequence &
Bidirectional global scan &
& \cmark & &
& & \cmark & \cmark \\

GLADMamba~\cite{fu2025gladmamba} &
Multi-view graph representations &
No central graph serialization &
& \cmark & &
& & \cmark & \\


% =========================================================
\midrule
\multicolumn{10}{l}{\textbf{Dynamic and Spatio-Temporal Graph Models}} \\
\midrule

GraphSSM~\cite{li2024state} &
Sequence of graph snapshots &
Temporal snapshot order &
& & \cmark &
\cmark & & & \\

GSSM~\cite{zhou2024graph} &
Base-station node time series &
Temporal order &
& \cmark & &
& \cmark & & \\

DG-Mamba~\cite{yuan2025dg} &
Successive dynamic graph states &
Selective scan across graph snapshots &
& \cmark & &
& & \cmark & \\

DyGMamba~\cite{ding2024dygmamba} &
Historical node interactions and temporal patterns &
Chronological interaction order &
\cmark & & &
\cmark & & & \\

DyG-Mamba~\cite{li2024dyg} &
Timestamped interaction sequences &
Irregular chronological event order &
& \cmark & &
\cmark & & & \\

STG-Mamba~\cite{li2024stg} &
Graph-aware node time-series representations &
Temporal processing across multiple granularities &
& \cmark & &
& & \cmark & \\

SpoT-Mamba~\cite{choi2024spot} &
Graph-derived neighborhood sequences &
BFS, DFS, random-walk, and temporal scanning &
\cmark & & &
\cmark & & & \\

PS-Mamba~\cite{dong2025ps} &
Human joints across pose sequences &
Four graph-guided bidirectional spatial--temporal scans &
\cmark & & &
& & \cmark & \\

FuzzMamba~\cite{chen2026modeling} &
Node-centric spatial neighborhoods with contextual features &
Adaptive spatial ordering varying over time &
& \cmark & &
\cmark & & & \\

STMAGRN~\cite{zhang2025spatio} &
Traffic sequences and dynamic graph representations &
Temporal sequence processing &
\cmark & & &
& \cmark & & \\


% =========================================================
\midrule
\multicolumn{10}{l}{\textbf{Heterogeneous and Higher-Order Graph Models}} \\
\midrule

HeteGraph-Mamba~\cite{pan2024hetegraph} &
Metapath-based heterogeneous graph tokens &
Hierarchical within-type and across-type processing &
\cmark & & &
\cmark & & & \\

Mamba-GTC~\cite{Meng2026Mamba-GTC} &
Multi-hop neighborhood tokens &
Hop-evolution ordering &
\cmark & & &
& & \cmark & \\

TopoMamba~\cite{montagna2024topological} &
Incident simplices aggregated by rank &
Rank-ordered sequence &
\cmark & & &
\cmark & & & \\

CCMamba~\cite{chen2026ccmamba} &
Multi-rank incidence representations &
Rank- and incidence-aware linearization &
\cmark & & &
\cmark & & & \\

\bottomrule

\end{tabular}%
}

\vspace{1mm}
\begin{minipage}{0.98\textwidth}
\footnotesize
\textit{Note:} D3 indicates where graph structure enters the state-space
computation: input-, dynamics-, or operator-level coupling.
D4 indicates how the state-space component is combined with graph-learning
mechanisms. Multiple check marks are used when an architecture combines
integration strategies.
\end{minipage}

\end{table*}


\subsection{Heterogeneous and Higher-Order Graph Models}
\label{subsec:heterogeneous-higherorder}

Heterogeneous and higher-order Graph Mamba models primarily address the challenge of preserving type, hop, rank, and incidence information during sequence construction and state-space processing. Their main architectural emphasis lies in how complex structural semantics are encoded through tokenization and ordering.

\paragraph{\textbf{Type-aware heterogeneous graph modeling.}}
HeteGraph-Mamba~\cite{pan2024hetegraph} introduces a hierarchical
tokenization strategy for heterogeneous graphs. For each target node,
metapath-based instances are used to construct graph tokens, while
representations from different node types are aligned in a common latent
space. The model then performs two levels of selective processing: dependencies
are first modeled among nodes of the same type and subsequently across
different node types. HeteGraph-Mamba therefore incorporates heterogeneous
semantics mainly through token construction (D1) and hierarchical type-aware
ordering (D2), allowing selective state-space processing to capture both
within-type and cross-type long-range dependencies.
% [REV-R13] Added (red): closing sentence answering the four design questions in order.
%In terms of the four design questions, HeteGraph-Mamba reads metapath-based tokens, orders them hierarchically within and across node types, couples the graph at the input level, and runs standalone.
Mamba-GTC~\cite{Meng2026Mamba-GTC} follows a different strategy by combining
local graph learning with ordered state-space processing. Its enhanced
Hop2Token mechanism converts multi-hop neighborhoods into hop-evolution
sequences, where representations are arranged according to their distance from
the target node. A GNN branch captures local, unordered topological
information, whereas a Mamba branch models the ordered hop-based sequence to
capture longer-range dependencies. The two representations are aligned through
a cross-view contrastive objective. Accordingly, Mamba-GTC combines
hop-structured tokenization and ordering (D1--D2) with explicit GNN--Mamba
integration (D4).
% [REV-R13] Added (red): closing sentence answering the four design questions in order.
%In terms of the four design questions, Mamba-GTC reads multi-hop neighborhood tokens, orders them by hop distance, couples the graph at the input level, and runs in parallel with a GNN branch aligned through contrastive learning.

\paragraph{\textbf{Higher-order state-space modeling.}}
TopoMamba~\cite{montagna2024topological} extends state-space processing from
pairwise graphs to simplicial complexes. For each target node, the model
collects the incident simplices and groups them according to their rank.
Representations belonging to the same rank are aggregated, and the resulting
rank-level representations form an ordered sequence processed by Mamba.
Information from higher-order cells of different ranks can therefore interact
within the same state-space computation without requiring a dedicated
higher-order message-passing rule. Its main architectural mechanisms are
higher-order token construction (D1) and rank-based serialization (D2).
% [REV-R13] Added (red): closing sentence answering the four design questions in order.
%In terms of the four design questions, TopoMamba reads incident simplices aggregated by rank, orders them by rank, couples the complex at the input level, and runs standalone.
CCMamba~\cite{chen2026ccmamba} generalizes this idea to combinatorial
complexes. Rather than constructing a sequence only around a target node,
CCMamba organizes multi-rank incidence relationships into structured,
rank-aware sequences and processes them using selective state-space models.
This reformulates higher-order message propagation as sequential state-space
processing, allowing information to propagate across cells of different ranks
and incidence relationships. CCMamba therefore places its main emphasis on
rank-aware representation construction (D1) and incidence-aware
linearization (D2), rather than on explicit graph-conditioned modification of
the selective SSM parameters.
% [REV-R13] Added (red): closing sentence answering the four design questions in order. [REV-NOTE] verify the coupling level / integration type against the original paper.
%In terms of the four design questions, CCMamba reads multi-rank incidence representations, orders them by rank and incidence, couples the complex at the input level, and runs standalone, replacing higher-order message passing.

Therefore, heterogeneous and higher-order Graph Mamba architectures emphasize tokenization and serialization more than direct modification of the state-space dynamics. Their main challenge is to preserve structural semantics such as type, hop, rank, and incidence information during sequence construction. Accordingly, D1 and D2 are particularly prominent in these settings.

% [REV-R16] Added (red): the dominant design question for heterogeneous and higher-order graphs.
%In terms of the four design questions, most structure in these settings is carried by tokenization, with ordering as its complement, while all four models keep coupling at the input level.


 
Table~\ref{tab:unified-taxonomy} consolidates the reviewed architectures under the four design dimensions. The comparison shows that most models incorporate graph structure at the input or dynamics level, while operator-level coupling remains relatively uncommon. Integration is dominated by standalone and parallel designs, whereas sequential and embedded configurations appear less frequently.
 
 
 

 % ===========================================================================
% 4.3  Domain-Specific Graph Mamba Applications  (v3: six data-type subsections)
% Needs in preamble: array, booktabs, colortbl, xcolor, graphicx,
%                    and \cmark \pmark \xmark (already in your main file).
% \rv{...} = red "new/changed" marker. \renewcommand{\rv}[1]{#1} switches it off.
% Figure: fig_design_patterns_v3.pdf, label fig:design-patterns (see figure source header).
% ===========================================================================
\providecommand{\rv}[1]{#1}
\newcolumntype{Y}[1]{>{\raggedright\arraybackslash}p{#1}}
\newcolumntype{Z}[1]{>{\centering\arraybackslash}p{#1}}

% Domain-specific applications section (compact, user-specified structure).
% Needs: array, booktabs, colortbl, xcolor, graphicx, \cmark-free; figure file fig_design_patterns_compact.pdf
\newcolumntype{Y}[1]{>{\raggedright\arraybackslash}p{#1}}
\newcolumntype{Z}[1]{>{\centering\arraybackslash}p{#1}}


 

 

%---------------------------------------------------------------------------



% ==========================================================================
% MERGED SECTION: old Section 4 (applications) + old Section 6 (comparison)
% - Model descriptions appear ONCE (first part of each subsection).
% - Each benchmark comparison is a closing \paragraph{Performance comparison.}
%   that points back to the model description with \ref and does not re-describe it.
% - Tables: keep tab:healthcare, tab:remote-visual, tab:other-areas exactly as in the
%   current Section 4; move the five comparison tables from old Section 6 to the
%   places marked ">>> MOVE TABLE HERE" (their code and labels do not change).
% - Labels used here that you must define/check:
%     sec:neuro, sec:clinical-bio, sec:histopath, sec:med-imaging  (subsubsections)
%     sec:st-graphs  (the subsection that introduces STG-Mamba / SpoT-Mamba, i.e. the
%                     dynamic and spatio-temporal graphs subsection and its Table 3)
% ==========================================================================
\section{Structural Properties and Computational Trade-offs}
\label{sec:structural-properties}

Building on the architectural taxonomy of Section~\ref{sec:taxonomy}, this
section examines the consequences of the design choices used to adapt selective
state-space models to graphs. We focus on three questions: how serialization affects permutation behavior,
how long-range dependencies are represented after graph-to-sequence
transformation, and whether the linear complexity of the selective scan
translates into end-to-end scalability. These aspects directly address RQ2.
% [REV-R18] Added (red): each subsection is tied to a design question.
%In terms of the four design questions, Section~\ref{subsec:order-sensitivity} concerns ordering, Section~\ref{subsec:long-range} concerns tokenization and ordering, Section~\ref{subsec:complexity} concerns the cost contributed by all four questions, and Section~\ref{subsec:design-implications} concerns the level of coupling.

%====================================================================
\subsection{Node Relabeling and Ordering Sensitivity}
\label{subsec:order-sensitivity}

% [REV-R18] Signpost added (red).
This subsection concerns the ordering question. Node relabeling and sequence ordering are related but distinct properties.
As defined in Section~\ref{sec:graph-foundations}, permutation equivariance
concerns whether a node-level model behaves consistently when the indices of
an otherwise identical graph are changed. Ordering sensitivity instead asks
whether a fixed graph produces different representations when its
graph-derived tokens are presented to the SSM in different admissible orders.
Let $\Pi(G)$ denote the set of orderings allowed by a serialization rule for a
fixed graph $G$. For an ordering $\pi\in\Pi(G)$, a serialization-based
state-space branch can be written abstractly as:

\begin{equation}
F_{\pi}(G)
=
P_{\pi}^{\top}
\mathcal{M}_{\mathrm{SSM}}
\left(P_{\pi}H\right),
\label{eq:order-output}
\end{equation}

where $H$ denotes the graph-derived representations and $P_{\pi}$ applies the
selected order before state-space processing. Independence from the chosen
serialization would require:

\begin{equation}
F_{\pi_1}(G)=F_{\pi_2}(G),
\qquad
\forall \pi_1,\pi_2\in\Pi(G).
\label{eq:order-independence}
\end{equation}

A conventional sequential selective SSM does not provide this property by
construction because its recurrent state depends on token position.
Consequently, the mechanism used in D2 to obtain a sequence can itself become
a source of model variability.
Graph-Mamba~\cite{wang2024graph} explicitly addresses this issue through
graph-centric node prioritization and permutation-based training. Nodes with
the same structural priority are randomly shuffled during training, while
multiple permuted outputs are averaged at inference. The authors motivate this
procedure as a way to reduce sequence-related bias and improve stability.
Their ablation study further shows that node-level permutation combined with
degree-based prioritization performs better than the alternative permutation
and prioritization strategies evaluated in the paper. 

GMN~\cite{behrouz2024graph} takes a different approach. Node or neighborhood
tokens are constructed and ordered before being processed with bidirectional
Mamba. Bidirectional processing removes the privileged causal direction of a
single forward scan, but the resulting computation still operates on a
particular token sequence.
GSSC~\cite{huang2024can} provides a contrasting design. Instead of serializing
nodes, it replaces the sequential convolution with permutation-equivariant
global aggregation based on graph positional information and factorizable
graph kernels. The original paper explicitly identifies permutation
equivariance as a property of the GSSC operator. In this case, permutation
behavior is addressed at the operator level rather than through training-time
randomization or repeated scans.

These designs illustrate an important distinction between an
\emph{architectural guarantee} and an \emph{empirical strategy for reducing
order dependence}. Permutation-equivariant operators such as GSSC provide the
former under their stated assumptions, whereas random permutation,
bidirectional processing, and alternative serialization strategies primarily
address the latter. 

%====================================================================
\subsection{Long-Range Dependency Modeling and State Retention}
\label{subsec:long-range}

% [REV-R18] Signpost added (red).
Long-range behavior depends jointly on tokenization and ordering, because positions in the processed sequence are produced by these two choices rather than by distance in the graph. Selective SSMs are attractive for graph learning partly because their latent state provides a mechanism for transmitting information across long sequences. However, after a graph is transformed into a sequence, distance in the state-space computation acquires a different meaning from distance in the original graph.

As shown by the retention coefficient $\alpha_{t\leftarrow k}$ defined in
Equation~\eqref{eq:alpha-retention}, the contribution of an earlier token
$u_k$ to a later output $y_t$ depends on the sequence of intermediate
state transitions. In serialization-based Graph Mamba models, however, the
positions $k$ and $t$ are determined by tokenization (D1) and ordering (D2),
rather than by shortest-path distance in the original graph.

This distinction is visible in Graph-Mamba~\cite{wang2024graph}. The model
places nodes according to graph-derived priority so that selected nodes receive
different amounts of preceding sequence context. Its long-range modeling
mechanism therefore depends jointly on the selective SSM and the imposed
ordering, rather than on the state-space recurrence alone.
In GMN~\cite{behrouz2024graph}, the unit processed by Mamba may itself contain
a sampled neighborhood. Consequently, long-range modeling occurs at two
levels: local graph structure is incorporated during neighborhood sampling and
encoding, while dependencies among the resulting tokens are subsequently
modeled by bidirectional Mamba. This also means that improvements cannot be
attributed to the state-space component independently of the tokenization
procedure.

GSSC~\cite{huang2024can} provides stronger operator-level evidence for
long-range graph interaction. Its authors prove that there exists a graph
kernel parameterization for which the gradient between two node
representations does not decay as their shortest-path distance increases.
Accordingly, the paper characterizes GSSC as both permutation-equivariant and
long-range without requiring node serialization.

For temporal graphs, the interpretation changes again. DyGMamba
~\cite{ding2024dygmamba} is explicitly designed to exploit extensive historical
interaction information in continuous-time dynamic graphs. Its node-level SSM
encodes historical interactions, while a time-level SSM captures temporal
patterns used to identify relevant historical information. Stronger performance is observed when substantial long-term temporal information is available, whereas limited historical context may reduce effectiveness.

Therefore, ``long range'' should not be used as a single undifferentiated
property in Graph Mamba. Depending on the architecture, it may refer to
distance between serialized tokens, shortest-path distance in the graph, or
elapsed time between interactions. Meaningful analysis should explicitly
distinguish these three notions:

\begin{equation}
\text{sequence distance}
\;\neq\;
\text{graph distance}
\;\neq\;
\text{temporal distance}.
\label{eq:distance-distinction}
\end{equation}

This distinction is particularly important when comparing architectures whose
tokenization and serialization strategies differ.

%====================================================================
\subsection{Computational Complexity and Scalability}
\label{subsec:complexity}

% [REV-R18] Signpost added (red).
Each design question contributes its own cost term, so the linear complexity of the selective scan describes only part of the pipeline. The selective scan underlying Mamba scales linearly with the length of the
processed sequence. This property is well established for the state-space
component, but it does not imply that every Graph Mamba architecture has
linear end-to-end cost. Graph-specific operations may precede, accompany, or
follow the scan.
To make this distinction explicit, we use the following conceptual
decomposition:

\begin{equation}
T_{\mathrm{total}}
=
T_{\mathrm{graph}}
+
T_{\mathrm{token}}
+
T_{\mathrm{serialize}}
+
T_{\mathrm{SSM}}
+
T_{\mathrm{fusion}}
+
T_{\mathrm{readout}},
\label{eq:end-to-end-cost}
\end{equation}

where the terms respectively denote graph construction or preprocessing,
token construction, ordering or traversal, selective state-space processing,
graph--SSM fusion, and task-specific readout. 
% [REV-R18] Added (red): mapping of the cost terms to the design questions.
Here, $T_{\mathrm{token}}$ reflects tokenization, $T_{\mathrm{serialize}}$ reflects ordering (including the number of scans), $T_{\mathrm{SSM}}$ together with any structural conditioning reflects coupling, and $T_{\mathrm{fusion}}$ reflects integration; $T_{\mathrm{graph}}$ belongs to graph construction, which we treat as context.

The original Graph-Mamba study~\cite{wang2024graph} reports linear
$O(L)$ complexity of its Graph-Mamba Block with respect to input sequence
length $L$, using Mamba's hardware-aware implementation. The paper also reports graphics processing unit (GPU) memory reductions
of up to 74\% in its large-graph experiments.
For larger graphs, however, the implementation introduces a binning strategy
that divides long node sequences into subsequences and applies prioritization
and permutation within each bin. Hence, its practical computational behavior
includes more than the selective scan alone.

GMN~\cite{behrouz2024graph} explicitly incorporates token construction into
its complexity analysis. For its subgraph-token setting, the reported
complexity is
$O\!\left(Ms(m+1)|V|+|E|\right),$
where $M$, $s$, and $m$ control neighborhood sampling and token construction.
When nodes themselves are used as tokens, the Mamba processing is linear in
$|V|$; with the optional message-passing component, the reported overall
dependence becomes $O(|V|+|E|)$. The paper consequently exposes an explicit
trade-off between richer/longer token sequences and computational cost.
GSSC~\cite{huang2024can} reports a complexity of
$O(nmd)$ for its graph state-space convolution, where $n$ is the number of
nodes and $m$ and $d$ denote the hidden and positional-encoding dimensions.
Importantly, the authors also evaluate preprocessing separately. When
Laplacian positional encodings are used, eigendecomposition can become a
non-negligible cost; for larger graphs, the paper therefore computes only the
required eigenpairs using iterative methods rather than relying on full
eigendecomposition.

Dynamic graph models introduce another form of scalability requirement because
the number of historical interactions can become large.
% [REV-R27] "45 GB" corrected to "48 GB" (NVIDIA A40 memory). [REV-NOTE] verify against the DyGMamba paper.
DyGMamba
~\cite{ding2024dygmamba}, for example, explicitly evaluates memory consumption
as the number of temporal neighbors increases. On the reported setup, the model
processes more than 8192 temporal neighbors on a 48\,GB NVIDIA A40 GPU, whereas
the compared DyGFormer configuration is reported to process up to 2048. This
result is specific to the experimental configuration of that paper.

Taken together, these studies show why \emph{scan complexity} and
\emph{end-to-end scalability} should be reported separately. The former is a
property of the state-space backbone, whereas the latter also depends on the
representation generated by D1, the ordering and number of scans in D2,
structural conditioning in D3, and additional graph branches in D4.
Accordingly, efficiency comparisons should not infer complete-pipeline
scalability from the asymptotic complexity of the selective scan alone.

%====================================================================
\subsection{Structural Trade-offs and Design Implications}
\label{subsec:design-implications}

The preceding analysis shows that the four taxonomy dimensions introduce
different benefits and costs rather than a single uniformly preferable Graph
Mamba design. Table~\ref{tab:structural-tradeoffs} summarizes these
relationships.

\begin{table*}[h]
\centering
\caption{
Structural and computational implications of the Graph Mamba taxonomy
dimensions.
}
\label{tab:structural-tradeoffs}
\small
\renewcommand{\arraystretch}{1.2}
\resizebox{\textwidth}{!}{%}
\begin{tabular}{@{}
p{1.0cm}
p{2.5cm}
p{4.5cm}
p{6.5cm}
@{}}
\toprule
\textbf{Dim.} &
\textbf{Design choice} &
\textbf{Structural implication} &
\textbf{Main trade-off} \\
\midrule

D1 &
% [REV-R18] Design-question name added (red).
\textbf{Tokenization:} Graph-derived token construction &
Controls which local, semantic, temporal, or higher-order information reaches
the SSM &
Richer tokens may improve structural context but introduce sampling,
encoding, or representation-construction cost \\

\addlinespace

D2 &
% [REV-R18] Design-question name added (red).
\textbf{Ordering:} Serialization and scan strategy &
Determines adjacency in the processed sequence and therefore the context seen
by the recurrent state &
Structurally meaningful order can improve context organization, but
serialization-based models may remain sensitive to the selected ordering \\

\addlinespace

D3 &
% [REV-R18] Design-question name added (red).
\textbf{Coupling:} Graph-conditioned state-space computation &
Allows graph, temporal, or structural information to influence state evolution
more directly &
Tighter graph--SSM coupling introduces additional architecture-specific
computation and reduces the ability to analyze the SSM independently \\

\addlinespace

D4 &
% [REV-R18] Design-question name added (red).
\textbf{Integration:} Graph--SSM integration and fusion &
Combines complementary local graph processing with state-space context
modeling &
Additional branches or fusion mechanisms increase end-to-end computation and
make performance gains harder to attribute to a single component \\

\bottomrule
\end{tabular}}
\end{table*}

To summarize, the taxonomy reveals three recurring design tensions. First, serialization provides a direct way to apply selective SSMs to graphs, but introduces an ordering choice that must be distinguished from permutation equivariance. Second, long-range state-space processing can connect distant sequence positions, yet the graph meaning of those positions depends on tokenization and ordering. Third, the linear complexity of the selective scan does not imply linear end-to-end graph processing, because graph construction, structural encoding, and auxiliary graph operations may add substantial cost. These tensions also reflect different coupling depths: input-level designs preserve a modular SSM but remain more exposed to ordering sensitivity, dynamics-level designs condition the state evolution on graph information at additional computational cost, and operator-level designs incorporate graph structure directly into the state-space operator, improving structural integration but reducing modularity. %Together, these observations provide the structural basis for the application analysis in Section~5 and the empirical evaluation that follows.

 

 
\section{Applications and Cross-Domain Design Patterns}
\label{sec:applications}
Graph Mamba has been applied across diverse domains in which graphs are often
constructed from the underlying data. This section examines these applications
and extends RQ1 by showing how the architectural choices identified earlier
recur in domain-specific settings.

Across the reviewed applications, four recurring design patterns emerge:
\textbf{dual-view (DV)}, where graph and SSM components process complementary
views or modalities; \textbf{relational--temporal (RT)}, where graph
representations capture relations and the SSM models temporal evolution;
\textbf{graph-guided sequence (GS)}, where graph structure guides tokenization,
ordering, or scanning; and \textbf{local-to-global (LG)}, where local graph
processing is combined with broader state-space context modeling. Each work is
assigned a dominant pattern, with a secondary pattern where appropriate.

These patterns reflect recurring combinations of the four design
questions introduced in Section~\ref{sec:unified-framework}. In RT designs,
temporal structure largely determines ordering; in GS designs, graph structure
guides tokenization or ordering; and DV and LG mainly differ in how graph and
state-space components are integrated. Figure~\ref{fig:design-patterns} summarizes the distribution of the four design patterns across application settings, while Table~\ref{tab:cross-domain-applications} provides a detailed cross-domain comparison of the reviewed applications and their corresponding graph--SSM patterns.

\begin{figure*}[!t]
\centering
\definecolor{pDV}{RGB}{176,74,123}
\definecolor{pRT}{RGB}{44,62,107}
\definecolor{pGS}{RGB}{197,138,18}
\definecolor{pLG}{RGB}{85,122,90}
\definecolor{dA}{RGB}{31,99,168}   % neural and physiological signals
\definecolor{dB}{RGB}{27,138,134}  % clinical, biological, molecular
\definecolor{dC}{RGB}{63,138,46}   % histopathology
\definecolor{dD}{RGB}{107,75,158}  % medical imaging
\definecolor{dE}{RGB}{217,122,24}  % remote sensing and visual
\definecolor{dF}{RGB}{79,93,117}   % other areas
% pattern tag + work name + citation (kept on one line)
\newcommand{\gw}[3]{\mbox{{\setlength{\fboxsep}{1.2pt}\colorbox{p#1}{\makebox[1.35em]{\color{white}\bfseries\tiny #1}}}~#2~\cite{#3}}}
\newcommand{\gwl}[2]{\mbox{{\setlength{\fboxsep}{1.2pt}\colorbox{p#1}{\makebox[1.35em]{\color{white}\bfseries\tiny #1}}}~#2}}
% panel header
\newcommand{\phd}[2]{{\color{#1}\bfseries\small #2}\\[4pt]}
% common panel width = widest header
\newlength{\pwA}\newlength{\pwB}
\settowidth{\pwA}{\sffamily\bfseries\small Neural and physiological signals}
\settowidth{\pwB}{\sffamily\bfseries\small Clinical, biological, molecular}
\ifdim\pwB>\pwA \setlength{\pwA}{\pwB}\fi
\settowidth{\pwB}{\sffamily\bfseries\small Remote sensing and visual data}
\ifdim\pwB>\pwA \setlength{\pwA}{\pwB}\fi
\addtolength{\pwA}{2pt}
\resizebox{\textwidth}{!}{%
\begin{tikzpicture}[font=\sffamily,
  panel/.style={rounded corners=6pt, draw=#1!55, line width=0.6pt, fill=#1!10,
                text width=\pwA, align=flush left, inner xsep=6pt, inner ysep=5pt,
                font=\sffamily\footnotesize},
  conn/.style={draw=#1, line width=1pt}]
% ---------- ring of data settings ----------
\def\ro{3.1}\def\ri{1.75}
\foreach \a/\b/\c in {30/90/dA,-30/30/dB,-90/-30/dC,210/270/dD,150/210/dE,90/150/dF}
  \fill[\c] (\a:\ri) arc (\a:\b:\ri) -- (\b:\ro) arc (\b:\a:\ro) -- cycle;
\foreach \a in {30,90,150,210,270,330} \draw[white,line width=4pt] (\a:1.6) -- (\a:3.25);
\fill[white] (0,0) circle (1.6);
\draw[black!25,line width=0.6pt] (0,0) circle (1.45);
\node[align=center,font=\sffamily\bfseries\normalsize] at (0,0.05) {Graph Mamba\\applications};
\foreach \m/\t in {60/{Neural\\signals},0/{Clinical \&\\molecular},-60/{Histo-\\pathology},240/{Medical\\imaging},180/{Remote \&\\visual},120/{Other\\areas}}
  \node[white,align=center,font=\sffamily\bfseries\scriptsize] at (\m:2.42) {\t};

% ---------- right column (anchored to middle panel) ----------
\node[panel=dB, anchor=west] (pB) at (3.6,0) {%
 \phd{dB}{Clinical, biological, molecular}
 \gw{DV}{MGDTA$^{\blacklozenge}$}{han2024innovative}\\
 \gw{DV}{MKHCNet$^{\blacklozenge}$}{lu2025mamba}\\
 \gw{RT}{MGSSM-SAKI$^{+\mathrm{DV}}$}{xu2024identifying}\\
 \gw{RT}{Aghaee et al.$^{\dagger}$}{aghaee2024graph}\\
 \gw{GS}{ExPath}{kotoge2026expath}};
\node[panel=dA, anchor=south west] (pA) at ([yshift=0.3cm]pB.north west) {%
 \phd{dA}{Neural and physiological signals}
 \gw{RT}{GraphS4mer$^{\dagger}$}{tang2023modeling}\\
 \gw{RT}{BrainMamba$^{+\mathrm{GS}}$}{behrouz2024brain}\\
 \gw{RT}{Brain-GM}{wang2024learning}\\
 \gw{RT}{MSGM}{liu2026msgm}\\
 \gw{LG}{Brain Network Mamba}{zhang2026brainnetworkmamba}};
\node[panel=dC, anchor=north west] (pC) at ([yshift=-0.3cm]pB.south west) {%
 \phd{dC}{Histopathology}
 \gw{DV}{MGCM$^{\blacklozenge}$}{cui2026mgcm}\\
 \gw{GS}{GraphMamba (WSI)}{zheng2025graphmamba}\\
 \gw{GS}{TopoMamSurv}{chen2026graph}\\
 \gw{LG}{GAT--Mamba}{ding2024combining}\\
 \gw{LG}{CGAM}{qu2025cgam}};

% ---------- left column (anchored to middle panel) ----------
\node[panel=dE, anchor=east] (pE) at (-3.6,0) {%
 \phd{dE}{Remote sensing and visual data}
 \gw{DV}{GraphMamba (HSI)}{yang2024graphmamba}\\
 \gw{DV}{MGF-GCN$^{\blacklozenge}$}{zhao2025mgf}\\
 \gw{RT}{Tang et al.$^{+\mathrm{GS}}$}{tang2025spatial}\\
 \gw{GS}{GraphMamba (tok.)}{ahmad2025graphmamba}\\
 \gw{GS}{GM-HAD}{tang2026graph}\\
 \gw{GS}{Hamba}{dong2024hamba}\\
 \gw{LG}{TGMN}{chu2025tgmn}};
\node[panel=dF, anchor=south east] (pF) at ([yshift=0.3cm]pE.north east) {%
 \phd{dF}{Other application areas}
 \gw{DV}{MambaForGCN}{lawan2024mambaforgcn}\\
 \gw{RT}{SAMBA}{mehrabian2024mamba}\\
 \gw{RT}{IDS--GraphMamba}{atitallah2025ids}\\
 \gw{RT}{Ren et al.}{ren2026spatio}};
\node[panel=dD, anchor=north east] (pD) at ([yshift=-0.3cm]pE.south east) {%
 \phd{dD}{Medical imaging}
 \gw{GS}{GGVMamba}{zhou2024efficient}\\
 \gw{LG}{GM-UNet}{zhang2024gm}\\
 \gw{LG}{HGM}{zhu2025hybrid}\\
 \gw{LG}{GMMN}{zhang2026graph}};

% ---------- connectors ----------
\foreach \ang/\p/\c in {60/pA/dA,0/pB/dB,-60/pC/dC}{
  \draw[conn=\c] (\ang:3.15) -- (\p.west);
  \fill[\c] (\p.west) circle (1.6pt);}
\foreach \ang/\p/\c in {120/pF/dF,180/pE/dE,240/pD/dD}{
  \draw[conn=\c] (\ang:3.15) -- (\p.east);
  \fill[\c] (\p.east) circle (1.6pt);}
\end{tikzpicture}}

\vspace{4pt}
{\sffamily\footnotesize
\begin{tabular}{@{}ll@{}}
\gwl{DV}{dual-view: graph and SSM process different views} & \gwl{RT}{relational--temporal: time orders the sequence}\\[2pt]
\gwl{GS}{graph-guided sequence: graph structure orders the sequence} & \gwl{LG}{local-to-global: local graph, then global SSM}
\end{tabular}\\[3pt]
$^{+X}$ secondary pattern;\quad $^{\blacklozenge}$ multimodal;\quad $^{\dagger}$ structured-SSM precursor without input-dependent selection.}

\caption{Landscape of domain-specific Graph Mamba applications. The central ring groups the surveyed works by data setting, and each callout panel lists the works of one setting. The coloured tag before each work marks its dominant graph--SSM design pattern: dual-view (DV), relational--temporal (RT), graph-guided sequence (GS), or local-to-global (LG). The pattern follows the data. Relational--temporal designs dominate time-indexed settings such as neural signals and network traffic. Local-to-global designs prevail in medical imaging. Graph-guided serialization is most common in remote sensing and histopathology, where no natural ordering of nodes exists. Superscripts denote secondary patterns, multimodal inputs, and structured-SSM precursors.}
\label{fig:design-patterns}
\end{figure*}


\subsection{Healthcare and Biomedical Applications}
\label{sec:healthcare}

Healthcare and biomedical applications provide a prominent example of these
cross-domain design patterns because they often combine relational structure
with temporal, spatial, or long-range dependencies. Depending on the task,
graph nodes may represent physiological variables, brain regions, clinical
features, molecular entities, tissue regions, or image-derived regions. The
reviewed works are organized into neural and physiological signals, clinical
and molecular data, histopathology, and medical imaging.


\paragraph{\textbf{Neural and physiological signals.}}
GraphS4mer~\cite{tang2023modeling} learns dynamic relationships among
physiological variables and combines Graph Structure Learning and Graph Isomorphism Network (GIN) with S4
layers for seizure detection, sleep staging, and electrocardiography (ECG)classification.
BrainMamba~\cite{behrouz2024brain} separates brain-network modeling through
BNMamba from temporal modeling through BTMamba, allowing spatial relations
among brain units and multivariate temporal dynamics to be processed
separately. Brain-GM~\cite{wang2024learning} models time-varying brain
connectivity for brain decoding and disease classification, while Brain
Network Mamba~\cite{zhang2026brainnetworkmamba} applies bidirectional
state-space processing to resting-state functional magnetic resonance imaging (fMRI) network representations.
MSGM~\cite{liu2026msgm} further combines multi-scale temporal segmentation
with local--global brain graphs for electroencephalography (EEG) emotion recognition. These approaches
mainly follow the relational--temporal pattern because graph structure
captures functional relationships while the state-space component models their
evolution over time.

\paragraph{\textbf{Clinical, biological, and molecular data.}}
MGSSM-SAKI~\cite{xu2024identifying} adaptively infers latent graphs from
multimodal clinical variables for subphenotyping sepsis with acute kidney
injury, combining graph-based local relationships with temporal patient-state
modeling. Aghaee et al.~\cite{aghaee2024graph} represent metabolic pathways as
graphs and use structured state-space modeling to capture their temporal
evolution. ExPath/PathMamba~\cite{kotoge2026expath} converts sampled molecular
pathways into sequences for Mamba while retaining graph-based pathway
structure. MGDTA~\cite{han2024innovative} combines drug molecular graphs with
Mamba-encoded protein sequences for drug--target affinity prediction, whereas
MKHCNet~\cite{lu2025mamba} integrates a disease semantic knowledge graph with
Mamba-based modeling of electronic health records for interpretable International Classification of Diseases (ICD) coding.
This group spans RT, graph-guided sequence, and dual-view
designs depending on whether graph and sequence information are temporally
coupled, serialized, or processed as complementary modalities.

\paragraph{\textbf{Histopathology.}}
Graph Mamba has also been adopted for Whole Slide Image (WSI) analysis, where
graphs typically represent spatial relationships among tissue regions.
GAT--Mamba~\cite{ding2024combining} combines graph attention with Mamba for
progression-free survival prediction in early-stage lung adenocarcinoma.
GraphMamba (WSI)~\cite{zheng2025graphmamba} introduces multi-level graph construction
with intra- and cross-group Graph Mamba modules for WSI classification.
MGCM~\cite{cui2026mgcm} combines transcriptomic co-expression and pathology
patch graphs with interactive Mamba modules for multimodal cancer survival
prediction. TopoMamSurv~\cite{chen2026graph} explicitly incorporates
topology-aware node ordering before bidirectional Mamba processing, whereas
CGAM~\cite{qu2025cgam} combines graph attention and Mamba for esophageal
pathology grading. These models predominantly follow local-to-global,
graph-guided sequence, or dual-view patterns depending on how local
tissue structure is coupled with broader state-space context.


\paragraph{\textbf{Medical imaging.}}
In medical imaging, graph construction is typically applied to image features
or anatomical regions before state-space processing. GM-UNet
~\cite{zhang2024gm} embeds dynamic feature graphs within a U-Net for medical
image segmentation. GGVMamba~\cite{zhou2024efficient} combines directed
scanning, a Graph Mamba encoder, and gender-adaptive graph structures for
pediatric bone-age assessment. HGM~\cite{zhu2025hybrid} combines graph
convolution with multi-directional Mamba for polyp segmentation, while
GMMN~\cite{zhang2026graph} integrates graph convolution and Mamba for macular
edema diagnosis from fundus images. Most of these models adopt a
local-to-global pattern in which graph operations capture local spatial or
topological relationships and Mamba provides broader contextual modeling.

 



The healthcare applications demonstrate that the graph is rarely a fixed
input independent of the task. Instead, it is commonly learned or constructed
from physiological connectivity, clinical variables, molecular relations, or
spatial tissue organization. Consequently, the application design is largely
determined by how graph construction is coupled with temporal modeling,
sequence formation, or multimodal fusion.

\subsection{Remote Sensing and Structured Visual Data}
\label{sec:remote-visual}

Remote sensing and structured visual applications often construct graphs from
spatial regions, superpixels, spectral--spatial neighborhoods, or articulated
joints. In these settings, graph processing preserves spatial or structural
relationships, while the state-space component models broader contextual,
spectral, or temporal dependencies. %Table~\ref{tab:remote-visual} summarizes
%the reviewed models according to graph construction, graph--SSM division of
%labor, and dominant design pattern.

\paragraph{\textbf{Remote sensing.}}
GraphMamba (HSI)~\cite{yang2024graphmamba} separates spectral and spatial
modeling by assigning high-dimensional spectral representations to HyperMamba
and spatial neighborhoods to SpatialGCN, forming a dual-view design.
MGF-GCN~\cite{zhao2025mgf} similarly combines height-aware graph convolution
with hierarchical cross-modal Mamba processing for multimodal remote-sensing
segmentation. Ahmad et al.~\cite{ahmad2025graphmamba} instead use graph
structure to construct and prioritize spectral--spatial tokens before Mamba
processing, corresponding to a graph-guided sequence design.
TGMN~\cite{chu2025tgmn} aggregates superpixel subgraphs using graph
convolution before applying a Mamba encoder for broader dependencies, following
a local-to-global pattern. GM-HAD~\cite{tang2026graph} constructs a
superpixel region-adjacency graph for hyperspectral anomaly detection and uses
the graph both for local aggregation and to guide the ordering of regions
processed by the selective SSM, making it a GS design.

\paragraph{\textbf{Articulated and human motion.}}
Graph-guided sequencing is particularly natural for articulated structures,
where the graph provides meaningful spatial relationships between joints.
Hamba~\cite{dong2024hamba} uses joint relations to determine a
Graph-guided Bidirectional Scan for 3D hand reconstruction, thereby aligning
the state-space scan with the hand's structural topology. Tang et
al.~\cite{tang2025spatial} model human skeletons as joint sequences within a
spatio-temporal Graph Mamba framework for music-guided dance synthesis,
combining graph-defined spatial relations with temporal state-space modeling.
These models therefore illustrate graph-guided sequence and relational--temporal patterns,
respectively.


% [REV-R19] Duplicate sentence removed (Table 7 is already introduced at the start of this subsection):
% [REV-R19-ORIGINAL] Table~\ref{tab:remote-visual} summarizes the reviewed remote-sensing and
% [REV-R19-ORIGINAL] structured-visual models in terms of task, graph construction, graph--SSM
% [REV-R19-ORIGINAL] division of labor, and dominant cross-domain design pattern.
The remote-sensing and structured-visual applications rely heavily on how
spatial structure is transformed into graph representations or scan
trajectories. Dual-view designs separate complementary modalities, local-to-global
models combine regional graph processing with broader state-space context, and
graph-guided sequence models use spatial topology directly to influence token
construction or scanning.

\subsection{Other Application Areas}
\label{sec:other-areas}

Beyond healthcare and visual applications, Graph Mamba has also been adopted
in financial forecasting, industrial prognostics, cybersecurity, and
language-based learning. In these settings, graphs typically represent
relationships among assets, sensors, communication flows, or words, while the
state-space component models temporal or contextual dependencies.
%Table~\ref{tab:other-areas} summarizes these applications according to graph construction, graph--SSM division of labor, and dominant design pattern.

\paragraph{\textbf{Financial forecasting.}}
SAMBA~\cite{mehrabian2024mamba} models financial markets by combining
Adaptive Graph Convolution over inter-asset relationships with Bidirectional
Mamba over historical price sequences. The graph captures dependencies among
assets, whereas the state-space component models temporal evolution, making the
architecture a relational--temporal design.


\paragraph{\textbf{Industrial prognostics.}}
Ren et al.~\cite{ren2026spatio} construct a spatio-temporal hypergraph from
multi-sensor degradation signals and prior engine knowledge for remaining
useful life prediction. The graph represents higher-order relations among
sensor and engine states, while Graph Mamba models their temporal evolution and
dependencies, corresponding primarily to a relational--temporal design.



\paragraph{\textbf{Cybersecurity and intrusion detection.}}
IDS--GraphMamba~\cite{atitallah2025ids} represents Internet of Medical Things (IoMT) communication flows
through a row-stochastic transition matrix and combines directional graph
convolution, Markov-based message passing, and a selective state-space module.
The graph component captures directional relationships among flows, whereas
the SSM models sequential and adaptive dependencies, following a relational--temporal pattern.

\paragraph{\textbf{Language and semantic modeling.}}
MambaForGCN~\cite{lawan2024mambaforgcn} combines a syntactic dependency graph
with Mamba-based contextual modeling for aspect-based sentiment analysis. The
graph convolutional branch captures structural relations among words, while
the Mamba branch models broader semantic and contextual dependencies before
fusion. Because the graph and SSM operate on complementary representations,
the model follows a dual-view design.





% [REV-R26] Placeholder section heading removed; Fig.~\ref{fig:design-patterns} is referenced from Section 5 and floats there.
% [REV-R26-ORIGINAL] \section{table and fig}


%---------------------------------------------------------------------------

\begin{landscape}
% ============================================================
% PAGE 1 — HEALTHCARE AND BIOMEDICAL APPLICATIONS
% ============================================================

\scriptsize
\setlength{\tabcolsep}{3pt}
%\renewcommand{\arraystretch}{1.03}

\begin{longtable}{@{}
p{2.2cm}
p{2.8cm}
p{3.8cm}
p{7cm}
cccc
@{}}

\caption{
Cross-domain Graph Mamba applications and recurring graph--SSM design patterns.
}
\label{tab:cross-domain-applications}\\

\toprule
\multirow{2}{*}{\textbf{Work}} &
\multirow{2}{*}{\textbf{Task}} &
\multirow{2}{*}{\textbf{Graph construction}} &
\multirow{2}{*}{\textbf{Graph--SSM division of labor}} &
\multicolumn{4}{c}{\textbf{Pattern}} \\

\cmidrule(lr){5-8}

& & & &
\rotatebox{90}{DV} &
\rotatebox{90}{RT} &
\rotatebox{90}{GS} &
\rotatebox{90}{LG} \\
\midrule
\endfirsthead

\toprule
\multirow{2}{*}{\textbf{Work}} &
\multirow{2}{*}{\textbf{Task}} &
\multirow{2}{*}{\textbf{Graph construction}} &
\multirow{2}{*}{\textbf{Graph--SSM division of labor}} &
\multicolumn{4}{c}{\textbf{Pattern}} \\

\cmidrule(lr){5-8}

& & & &
\rotatebox{90}{DV} &
\rotatebox{90}{RT} &
\rotatebox{90}{GS} &
\rotatebox{90}{LG} \\
\midrule
\endhead

\bottomrule
\endlastfoot


\multicolumn{8}{l}{\textbf{Healthcare and Biomedical Applications}}\\
\midrule

\rowcolor{gray!12}
\multicolumn{8}{l}{\textbf{\textit{Neural and physiological signals}}} \\

GraphS4mer$^{\dagger}$~\cite{tang2023modeling} &
Seizure, sleep, ECG classification &
Physiological channels; dynamically learned edges &
GIN captures graph relations; S4 models long-range temporal dependencies. &
& \cmark & & \\

BrainMamba~\cite{behrouz2024brain} &
Brain encoding; attention-deficit/hyperactivity disorder and seizure detection &
Brain-unit tokens and multivariate time series &
Message passing models local relations; selective SSM captures node and temporal dependencies. &
& \cmark & \cmark & \\

Brain-GM~\cite{wang2024learning} &
Brain decoding and disease classification &
Dynamic brain networks &
Graph learning models time-varying connectivity; SSM captures temporal evolution. &
& \cmark & & \\

Brain Network Mamba~\cite{zhang2026brainnetworkmamba} &
Resting-state fMRI analysis &
Functional brain networks &
Graph structure captures regional interactions; bidirectional SSM models broader dependencies. &
& & & \cmark \\

MSGM~\cite{liu2026msgm} &
EEG emotion recognition &
Local--global brain graphs &
Graph modules model brain relations; SSM captures spatio-temporal dependencies. &
& \cmark & & \\


\midrule
\rowcolor{gray!12}
\multicolumn{8}{l}{\textbf{\textit{Clinical, biological, and molecular data}}} \\

MGSSM-SAKI~\cite{xu2024identifying} &
Sepsis-AKI subphenotyping &
Adaptive graph from multimodal clinical variables &
Graph aggregation models patient-state relations; SSM captures temporal evolution. &
\cmark & \cmark & & \\

Aghaee et al.$^{\dagger}$~\cite{aghaee2024graph} &
Metabolic-pathway modeling &
Metabolic-pathway graphs &
Graph models pathway interactions; structured SSM captures temporal evolution. &
& \cmark & & \\

ExPath / PathMamba~\cite{kotoge2026expath} &
Pathway inference and explanation &
Molecular network and sampled pathway sequences &
GIN captures local pathway structure; Mamba processes graph-derived pathways. &
& & \cmark & \\

MGDTA~\cite{han2024innovative} &
Drug--target affinity prediction &
Drug graph and protein sequence &
Graph encoder models molecular structure; Mamba models protein sequence before fusion. &
\cmark & & & \\

MKHCNet~\cite{lu2025mamba} &
Automatic ICD coding &
Disease knowledge graph linked to ICD labels &
Graph enriches label semantics; Mamba captures long-range EHR-text dependencies. &
\cmark & & & \\


\midrule
\rowcolor{gray!12}
\multicolumn{8}{l}{\textbf{\textit{Histopathology}}} \\

GAT--Mamba~\cite{ding2024combining} &
Lung-cancer survival prediction &
WSI tissue-region graph &
GAT models local tissue relations; Mamba captures global contextual dependencies. &
& & & \cmark \\

GraphMamba (WSI)~\cite{zheng2025graphmamba} &
WSI classification &
Multi-level instance/group graphs &
Graph hierarchy defines groups; Mamba models intra- and cross-group dependencies. &
& & \cmark & \\

MGCM~\cite{cui2026mgcm} &
Multimodal cancer survival prediction &
Transcriptomic and pathology graphs &
Graph convolution models within-modality structure; Mamba captures cross-modal dependencies. &
\cmark & & & \\

TopoMamSurv~\cite{chen2026graph} &
WSI survival analysis &
Topology-aware WSI graph &
Graph convolution models local structure; bidirectional Mamba processes topology-ordered nodes. &
& & \cmark & \\

CGAM~\cite{qu2025cgam} &
Esophageal pathology grading &
Tumor and microenvironment graph &
Graph attention models local topology; Mamba captures broader dependencies. &
& & & \cmark \\


\midrule
\rowcolor{gray!12}
\multicolumn{8}{l}{\textbf{\textit{Medical imaging}}} \\

GM-UNet~\cite{zhang2024gm} &
Image segmentation &
Dynamic feature graph within U-Net &
Graph modules model feature relations; Mamba provides broader contextual modeling. &
& & & \cmark \\

GGVMamba~\cite{zhou2024efficient} &
Bone-age assessment &
Gender-adaptive inter-region graph &
Graph models inter-region relations; graph-guided bidirectional Mamba captures context. &
& & \cmark & \\

HGM~\cite{zhu2025hybrid} &
Polyp segmentation &
Graph over image regions &
GCN captures local topology; multi-directional Mamba models global context. &
& & & \cmark \\

GMMN~\cite{zhang2026graph} &
Macular edema diagnosis &
Graph over fundus representations &
GCN captures local spatial correlations; Mamba models long-range dependencies. &
& & & \cmark \\

%%%%%%===============================================
\midrule
\multicolumn{8}{l}{\textbf{Remote Sensing and Structured Visual Applications}}\\
\midrule

\rowcolor{gray!12}
\multicolumn{8}{l}{\textbf{\textit{Remote sensing}}} \\

GraphMamba (HSI)~\cite{yang2024graphmamba} &
Hyperspectral classification &
Spatial neighborhoods and spectral representations &
SpatialGCN captures local spatial relations; HyperMamba models global spectral dependencies. &
\cmark & & & \\

MGF-GCN~\cite{zhao2025mgf} &
Multimodal remote-sensing segmentation &
Height-aware multimodal graph &
Graph convolution captures spatial/elevation relations; Mamba models cross-modal interactions. &
\cmark & & & \\

GraphMamba (tok.)~\cite{ahmad2025graphmamba} &
Hyperspectral classification &
Graph-guided spectral--spatial tokens &
Graph constructs and prioritizes tokens; Mamba models their global dependencies. &
& & \cmark & \\

TGMN~\cite{chu2025tgmn} &
Hyperspectral classification &
Superpixel subgraphs &
GCN models local superpixel relations; Mamba captures broader dependencies. &
& & & \cmark \\

GM-HAD~\cite{tang2026graph} &
Hyperspectral anomaly detection &
Superpixel region-adjacency graph &
Graph guides local aggregation and ordering; SSM processes the resulting sequence. &
& & \cmark & \\


\midrule
\rowcolor{gray!12}
\multicolumn{8}{l}{\textbf{\textit{Articulated and human motion}}} \\

Hamba~\cite{dong2024hamba} &
3D hand reconstruction &
Hand-joint graph &
Graph structure determines the bidirectional scan; SSM models the joint sequence. &
& & \cmark & \\

Tang et al.~\cite{tang2025spatial} &
Music-guided dance synthesis &
Dynamic human-skeleton graph &
Graph models joint relations; Graph Mamba captures spatial and temporal dependencies. &
& \cmark & \cmark & \\


\midrule
\multicolumn{8}{l}{\textbf{Other Application Areas}}\\
\midrule

SAMBA~\cite{mehrabian2024mamba} &
Stock-return forecasting &
Adaptive inter-asset graph &
Graph convolution models asset relations; bidirectional Mamba captures price dynamics. &
& \cmark & & \\

Ren et al.~\cite{ren2026spatio} &
Aero-engine remaining useful life (RUL)prediction &
Spatio-temporal sensor hypergraph &
Hypergraph models higher-order sensor relations; Graph Mamba captures temporal degradation. &
& \cmark & & \\

IDS--GraphMamba~\cite{atitallah2025ids} &
IoMT intrusion detection &
Flow transition graph &
Graph modules capture directional flow relations; SSM models sequential dependencies. &
& \cmark & & \\

MambaForGCN~\cite{lawan2024mambaforgcn} &
Aspect-based sentiment analysis &
Syntactic dependency graph &
SynGCN models syntax; Mamba models broader semantic context before fusion. &
\cmark & & & \\

\end{longtable}

\vspace{-2mm}

\begin{center}
\footnotesize
\textbf{Patterns}: DV = dual-view; RT = relational--temporal;
GS = graph-guided sequence; LG = local-to-global.
Multiple check marks indicate combined patterns.
$^{\dagger}$ denotes an S4-type or structured SSM precursor without
Mamba-style input-dependent selection.
\end{center}

\end{landscape}

 
\subsection{Cross-Domain Design Patterns}
\label{sec:cross-domain-patterns}

Across application domains, Graph Mamba models exhibit recurring ways of
combining graph structure with state-space processing. These patterns are
largely determined by the structure of the underlying data rather than by the
application domain itself.
The {relational--temporal} pattern is most common in settings where
relationships evolve over time, including neural signals, clinical data,
financial markets, industrial systems, and network traffic. In these models,
the graph represents interactions among entities, while the SSM captures their
temporal evolution.
The {dual-view} pattern appears mainly in multimodal settings where
graph and sequential representations provide complementary information, such
as molecular--protein modeling, multimodal pathology, remote sensing, and
language applications. Here, the graph and SSM process different views before
their representations are fused.
The {graph-guided sequence} pattern is used when graph structure
directly influences token construction, ordering, or scanning. This is common
in structured visual data, hyperspectral imaging, pathway modeling, and
topology-aware WSI analysis. Because the resulting sequence depends on graph
structure, these models are particularly sensitive to serialization choices.
Finally, the {local-to-global} pattern combines local graph
processing with broader state-space modeling. It is frequently used in
histopathology and medical imaging, where graph operations capture local
spatial relationships and Mamba provides wider contextual modeling.
A common observation across all four patterns is that the graph is usually
constructed rather than naturally given. Consequently, reported performance
reflects not only the state-space architecture, but also graph construction,
tokenization, ordering, and fusion choices.
% [REV-R19] Added (red): cross-domain finding stated with the design questions. [REV-NOTE] verify against Tables 6--8.
%In terms of the four design questions, graph structure in these applications enters almost exclusively at the input level, through graph construction, tokens, ordering, or a separate graph branch, rather than by modifying the selective state dynamics; the application-specific answers therefore lie mainly in the ordering and integration questions. 


\section{Comparative Evidence and Performance Evaluation}
\label{sec:evaluation}
This section addresses RQ3 by examining under which graph settings, tasks,
datasets, and evaluation conditions the reported advantages of Graph Mamba are
supported. We synthesize empirical evidence across the application settings
discussed in Section~\ref{sec:applications}, while accounting for differences
in datasets, graph construction, preprocessing, serialization strategies, and
experimental protocols.

\subsection{Evaluation Protocol and Evidence Provenance}
\label{sec:evidence-protocol}

Graph Mamba models are evaluated under different datasets, data splits,
preprocessing procedures, graph construction strategies, serialization
schemes, positional encodings, and training settings. Consequently, numerical
results reported across different studies should not automatically be
interpreted as controlled head-to-head comparisons.

To make the subsequent comparisons more transparent, we distinguish between
two levels of evidence. \textbf{Comparable published results} refer to results
reported under sufficiently similar datasets, tasks, metrics, and evaluation
protocols to support a meaningful comparison. \textbf{Contextual results}
refer to results obtained under substantially different settings and are
therefore used only to indicate broader performance trends. Unless explicitly
stated otherwise, all numerical results reported in this survey are taken from
the original publications and were not reproduced by us.

Accordingly, the performance tables in this section should be interpreted
together with their evaluation context. Differences in graph construction,
tokenization, serialization, training protocol, and implementation may
contribute to the reported performance in addition to the state-space
architecture itself.

\subsection{Benchmark Datasets and Tasks}
\label{sec:benchmarks}

The surveyed Graph Mamba models are evaluated across heterogeneous benchmark
settings, ranging from native graph datasets to graphs constructed from
biomedical signals, images, transportation networks, financial assets, and
language data. Table~\ref{tab:benchmark-summary} summarizes the main benchmarks
used in the comparative analyses of this survey, together with their task,
graph setting, evaluation metrics, and representative Graph Mamba models.

The diversity of these benchmarks highlights an important evaluation
challenge: identical metrics do not necessarily imply directly comparable
experimental conditions. In particular, graph construction, sequence
formation, preprocessing, and data splits differ across domains and
studies. The benchmark summary is therefore used primarily to establish the
evaluation context for the performance comparisons that follow.


\begin{table*}[h]
\centering
\caption{
Representative benchmarks used to evaluate Graph Mamba models.
}
\label{tab:benchmark-summary}
\scriptsize
\setlength{\tabcolsep}{4pt}
\renewcommand{\arraystretch}{1.2}
\resizebox{\textwidth}{!}{%
\begin{tabular}{Y{2.2cm} Y{2.8cm} Y{3.2cm} Y{2.2cm} Y{2.2cm}}
\toprule
\textbf{Benchmark} &
\textbf{Domain / Task} &
\textbf{Graph Setting} &
\textbf{Main Metrics} &
\textbf{Representative Models} \\
\midrule

Peptides-func, Peptides-struct &
General graph learning &
Static molecular graphs &
AP, MAE &
Graph-Mamba~\cite{wang2024graph}, GMN~\cite{behrouz2024graph}, GSSC~\cite{huang2024can} \\
\hline

BVFC-MEG, HCP-Mental, HCP-Age &
Brain classification &
Brain-network / physiological graphs &
Accuracy &
GraphS4mer~\cite{tang2023modeling}, BrainMamba~\cite{behrouz2024brain} \\
\hline

Indian Pines, Salinas, UH2013 &
Hyperspectral image classification &
Spectral--spatial graphs &
OA, AA, Kappa &
GraphMamba (HSI)~\cite{yang2024graphmamba} \\
\hline

PeMS04, HZMetro, KnowAir &
Spatio-temporal forecasting &
Traffic / environmental graphs &
RMSE, MAE, MAPE &
STG-Mamba~\cite{li2024stg} \\
\hline

NASDAQ, NYSE, DJIA &
Financial forecasting &
Inter-asset graphs &
RMSE, training time &
SAMBA~\cite{mehrabian2024mamba} \\
\hline

Restaurant14, Laptop14, Twitter14 &
Aspect-based sentiment analysis &
Syntactic dependency graphs &
Accuracy, F1-score &
MambaForGCN~\cite{lawan2024mambaforgcn} \\

\bottomrule
\end{tabular}}
{\footnotesize
\textbf{Metrics:} AP = Average Precision; MAE = Mean Absolute Error;
RMSE = Root Mean Squared Error; MAPE = Mean Absolute Percentage Error;
OA = Overall Accuracy; AA = Average Accuracy; Kappa = Cohen's Kappa coefficient.
}

\end{table*}


\subsection{Performance Comparison}
\label{sec:performance-comparison}

The available comparisons span different domains, datasets, and evaluation
protocols. Accordingly, the results are interpreted using the evidence
categories defined in Section~\ref{sec:evidence-protocol}. Rather than
establishing a universal ranking of Graph Mamba models, this subsection
summarizes where reported gains are supported under sufficiently comparable
benchmark settings. Detailed numerical comparisons are provided in
Appendix~\ref{app:performance-tables}.

\paragraph{\textbf{Brain and physiological data.}}
On BVFC-MEG, HCP-Mental, and HCP-Age, BrainMamba~\cite{behrouz2024brain}
reports higher mean accuracy than GraphS4mer~\cite{tang2023modeling}, with the
clearest advantage observed on BVFC-MEG
(Table~\ref{tab:Healthcompare}). However, the gains are less consistent on
HCP-Mental and HCP-Age: GraphS4mer falls below several non-SSM baselines, and
BrainMamba's margin over BrainMixer~\cite{behrouz2023learning} is small relative
to the reported variability. These results therefore provide stronger evidence
for selective state-space modeling in high-temporal-resolution settings than for
a general advantage across brain benchmarks~\cite{behrouz2024brain}. Both models follow the relational–temporal pattern. They differ mainly in the temporal backbone, since BrainMamba uses input-dependent selection whereas GraphS4mer relies on the non-selective S4 formulation.

\paragraph{\textbf{Remote sensing.}}
On Indian Pines, Salinas, and UH2013, GraphMamba (HSI)~\cite{yang2024graphmamba}
reports higher OA, AA, and Kappa than the compared recurrent, convolutional,
and Transformer-based baselines (Table~\ref{tab:RScomparison}). Its dual-view
design combines SpatialGCN for local spatial relationships with HyperMamba for
spectral dependency modeling, supporting the benefit of complementary spatial
and spectral representations under the reported setting.

\paragraph{\textbf{Spatio-temporal forecasting.}}
Across PeMS04, HZMetro, and KnowAir, STG-Mamba reports the
lowest RMSE and MAE among the compared methods
(Table~\ref{tab:TRperformance_comparison}). STG-Mamba~\cite{li2024stg}
combines graph-based spatial processing through KFGN with selective temporal
state evolution, corresponding to dynamics-level coupling and embedded
integration. Relative to the strongest baseline, the attention-based
STAEformer, STG-Mamba reduces RMSE by approximately 2.2\%, 2.4\%, and 7.4\%
on PeMS04, HZMetro, and KnowAir, respectively. Its MAPE is slightly higher on
PeMS04 (12.11\% vs. 11.98\%), while the LSTM-based baseline reports the lowest
MAPE on HZMetro (9.90\%). Because STG-Mamba and STAEformer differ in both
their spatial and temporal components, and no uncertainty estimates or
component-level ablations are reported, these results support the
competitiveness of the overall design but do not isolate the contribution of
the selective SSM~\cite{li2024stg}.

\paragraph{\textbf{Financial forecasting.}}
On NASDAQ, NYSE, and DJIA, SAMBA~\cite{mehrabian2024mamba} reports lower RMSE
than the LSTM~\cite{moghar2020stock}, Transformer~\cite{wang2022stock}, and
FourierGNN~\cite{yi2024fouriergnn} baselines considered in the original
comparison (Table~\ref{tab:FMcomparison}). Its reported training time is also
lower than those of the Transformer and FourierGNN, although higher than that
of the LSTM. The results therefore indicate a favorable
predictive--computational trade-off within the reported setting
~\cite{mehrabian2024mamba}.

\paragraph{\textbf{Language applications.}}
For aspect-based sentiment analysis, MambaForGCN~\cite{lawan2024mambaforgcn}
reports higher accuracy and F1-score than the compared graph-based baselines on
Restaurant14, Laptop14, and Twitter14
(Table~\ref{tab:SAperformance}). These results support the complementary use of
syntactic graph structure and state-space contextual modeling under the reported
evaluation protocol~\cite{lawan2024mambaforgcn}.

Taken together, the published comparisons show competitive gains in several
domain-specific settings, but the strength of the evidence varies across
datasets, metrics, baselines, and experimental protocols. The reported results
should therefore be interpreted as task-specific evidence rather than as a
general cross-domain superiority claim.

\subsection{Efficiency and Scalability}
\label{sec:efficiency-eval}

Beyond predictive performance, Graph Mamba is frequently motivated by the
linear sequence-length complexity of selective state-space models. However,
end-to-end efficiency also depends on graph construction, tokenization,
serialization, local graph processing, and fusion operations. Reported
efficiency should therefore be interpreted at the level of the complete
pipeline rather than the selective scan alone.

Several studies provide model-specific evidence of favorable scaling.
Graph-Mamba~\cite{wang2024graph} reports linear complexity for its Mamba-based
global block and reduced GPU memory usage on large graphs, although additional
costs arise from node prioritization and sequence construction. GMN
~\cite{behrouz2024graph} similarly maintains linear dependence on the number of
processed graph tokens, with the total cost additionally determined by
neighborhood sampling and optional message passing. GSSC~\cite{huang2024can}
reports linear-time state-space propagation with respect to graph size for its
structured kernel formulation, while preprocessing such as positional or
spectral encodings may introduce additional cost. 
For dynamic graphs, DyGMamba~\cite{ding2024dygmamba} provides empirical
scalability evidence for long interaction histories, supporting substantially
larger neighborhoods than attention-based dynamic-graph baselines under the
reported hardware configuration. These results illustrate an important
distinction: the state-space component may scale linearly with sequence length,
but practical scalability depends on the full graph-to-sequence pipeline.

Accordingly, meaningful efficiency comparisons should report not only
asymptotic complexity, but also parameter count, preprocessing cost, peak
memory, training and inference time, graph size, sequence length, and hardware
configuration.

\begin{table*}[h]
\centering
\caption{
Representative efficiency and scalability evidence reported for Graph Mamba models.
}
\label{tab:efficiency-summary}
\scriptsize
\setlength{\tabcolsep}{4pt}
\renewcommand{\arraystretch}{1.2}
\resizebox{\textwidth}{!}{%
\begin{tabular}{Y{2.4cm} Y{3.0cm} Y{4.1cm} Y{5.2cm}}
\toprule
\textbf{Model} &
\textbf{Reported scaling} &
\textbf{Additional cost sources} &
\textbf{Reported evidence} \\
\midrule

Graph-Mamba~\cite{wang2024graph} &
Linear Mamba global block &
Node prioritization, permutation, sequence construction &
Reports reduced GPU memory usage on large-graph experiments \\

GMN~\cite{behrouz2024graph} &
Linear in processed graph-token sequence length &
Neighborhood sampling and optional message passing &
Complexity depends on sampling parameters and graph size \\

GSSC~\cite{huang2024can} &
Linear-time structured state-space propagation &
Graph-kernel and positional/spectral preprocessing &
Reports graph-size-dependent linear propagation complexity \\

DyGMamba~\cite{ding2024dygmamba} &
Scales to long temporal interaction histories &
Historical-neighbor retrieval and temporal sequence construction &
Reports substantially larger supported neighborhoods than the compared
attention-based dynamic-graph baseline under the stated hardware setting \\

\bottomrule
\end{tabular}}
\end{table*}

\subsection{Reproducibility and Evaluation Gaps}
\label{sec:reproducibility}

Reproducibility remains difficult across Graph Mamba studies because evaluation
protocols differ in graph construction, preprocessing, serialization, data
splits, positional encodings, model configuration, and hardware. Even when the
same benchmark is used, these differences can substantially affect the input
presented to the state-space component and therefore complicate direct
comparison.

A further limitation is the incomplete reporting of computational settings.
Many studies emphasize the linear complexity of Mamba without consistently
reporting preprocessing cost, peak memory, sequence length, graph size, number
of scans, parameter count, or inference latency. Multi-seed evaluation is also
not uniformly adopted, making it difficult to distinguish systematic gains
from variation due to initialization, sampling, or graph ordering.

Ordering sensitivity is particularly under-evaluated. Models based on
graph-guided serialization may depend strongly on how nodes, regions, or paths
are arranged before state-space processing, yet controlled comparisons across
alternative orderings are still uncommon. Similarly, few studies isolate the
individual contributions of graph construction, local graph propagation,
serialization, and the selective SSM through complete ablation studies.

Future evaluations would therefore benefit from standardized reporting of data
splits, graph construction, serialization rules, random seeds, computational
resources, parameter counts, memory usage, and runtime. Controlled ablations
should further separate the effects of graph processing and state-space
modeling. Such practices would make comparisons across Graph Mamba
architectures more reliable and clarify which gains arise specifically from
selective state-space processing.

Overall, current evidence shows that Graph Mamba can achieve competitive
predictive performance and favorable scaling across several domains, but the
strength of the conclusions is limited by heterogeneous evaluation protocols
and incomplete efficiency reporting. More standardized and reproducible
experiments are therefore necessary before broad claims about accuracy,
scalability, or long-range modeling can be established across graph-learning
settings.

\section{Challenges and Research Agenda}
\label{sec6}

This section addresses RQ4 by linking the main open challenges to the four design questions. Ordering robustness relates to D2, expressive power to D1 and D3, efficient and hybrid designs to D3 and D4, and transferable pre-training to D1 and D2.

\subsection{Serialization, Permutation, and Robustness}
\label{sec:ch-serialization-robustness}

Order sensitivity remains a central limitation of serialization-based Graph
Mamba architectures. Different valid serializations of the same graph can
produce different state trajectories and predictions, even when the underlying
topology is unchanged. Existing approaches mitigate this dependence through
structure-aware prioritization~\cite{wang2024graph}, neighborhood
tokenization~\cite{behrouz2024graph}, spectral or temporal
ordering~\cite{zhao2024grassnet,choi2024spot}, bidirectional
scanning~\cite{behrouz2024graph,dong2024hamba}, and multi-order aggregation.
These strategies reduce order sensitivity but generally add computation and do
not provide a universal permutation guarantee.

For a graph $G$, let $\Pi(G)$ denote the set of admissible serializations and
let
$F_{\pi}(G)=
\mathbf{P}_{\pi}^{\top}
\mathcal{M}_{\mathrm{SSM}}
(\mathbf{P}_{\pi}\mathbf{H})$
denote the restored output under ordering $\pi$. Serialization sensitivity can
be measured as

\begin{equation}
\label{eq:order-sensitivity}
\mathcal{S}(G)
=
\mathbb{E}_{\pi,\pi'\sim\Pi(G)}
\left[
\frac{
\lVert F_{\pi}(G)-F_{\pi'}(G)\rVert_F
}{
\lVert F_{\pi}(G)\rVert_F
}
\right].
\end{equation}

Permutation-equivariant operators such as GSSC~\cite{huang2024can} satisfy
$\mathcal{S}(G)=0$ by construction, whereas serialization-based models may
exhibit non-zero sensitivity. Future evaluations should therefore report
performance across multiple admissible orderings, including variability and
prediction changes induced solely by re-ordering.

Robustness to structural and feature perturbations is also underexplored.
DG-Mamba~\cite{yuan2025dg} explicitly evaluates structural and feature
perturbations as well as evasion and poisoning attacks, while
DyG-Mamba~\cite{li2024dyg} uses spectral constraints to limit the influence of
noisy historical interactions. However, these mechanisms do not isolate the
robustness of the selective SSM itself. Moreover, structure-based serializers
introduce an additional attack surface: small changes to degree, centrality,
or traversal structure may reorder large parts of the sequence, allowing a
local graph perturbation to propagate non-locally through the state-space
recurrence.

Distribution shift raises related concerns. Changes in graph size become
changes in sequence length, while shifts in degree distribution, community
structure, or temporal dynamics alter the statistics of the generated
serialization. Existing evaluations rely predominantly on in-distribution
splits, with limited systematic analysis under structural, size, or temporal
shift~\cite{li2025ood,gui2022good}.

The main challenge is therefore to develop graph-aware state-space models that
combine low serialization sensitivity, permutation consistency, and robustness
without costly multi-order processing. 


\subsection{Expressive Power and Theoretical Understanding}
\label{sec:ch-expressivity}

% [REV-R21] Signpost added (red): design question addressed by this subsection.
Expressive power depends on how the tokenization and coupling questions are answered, rather than on the selective scan alone.
Message-passing GNNs are at most as powerful as the 1-dimensional
Weisfeiler--Lehman (1-WL) test in distinguishing non-isomorphic
graphs~\cite{xu2018powerful,morris2019weisfeiler}. Graph Transformers owe their
power largely to positional and structural encodings: with sufficiently
expressive encodings they are universal approximators, whereas without them
global attention cannot separate graphs that 1-WL cannot
separate~\cite{kreuzer2021rethinking,muller2023attending}. Where Graph Mamba
stands relative to these references is only partially understood;
Table~\ref{tab:expressivity} collects the available results.

\begin{table*}[!htbp]
\centering
\caption{Theoretical expressive-power results reported for
SSM-based graph models, with message-passing GNNs and graph Transformers as
references. PE/SE: positional/structural encodings.}
\label{tab:expressivity}
\small
\renewcommand{\arraystretch}{1.2}
\resizebox{\textwidth}{!}{%
\begin{tabular}{>{\raggedright\arraybackslash}p{3.0cm} >{\raggedright\arraybackslash}p{2.8cm} >{\raggedright\arraybackslash}p{5.6cm} >{\raggedright\arraybackslash}p{5.0cm}}
\toprule
\textbf{Model} &
\textbf{Permutation property} &
\textbf{Reported result} &
\textbf{Caveat} \\
\midrule

MPNNs (reference)~\cite{xu2018powerful,morris2019weisfeiler} &
Equivariant &
Upper-bounded by 1-WL &
Cannot count cycles; limited long-range propagation \\

\midrule

Graph Transformers (reference)~\cite{kreuzer2021rethinking,muller2023attending} &
Equivariant given equivariant PE/SE &
Universal approximation with sufficiently expressive PE/SE &
Power resides in the PE/SE; quadratic cost of global attention \\

\midrule

GSSC~\cite{huang2024can} &
Equivariant by construction &
Strictly more powerful than 1-WL and not more powerful than 3-WL; counts
3-paths and 3-cycles, and 4-paths and 4-cycles when the selection mechanism is
added &
Not a serialization-based model; relies on a factorizable distance kernel
built from PE \\

\midrule

GMN~\cite{behrouz2024graph} &
Not guaranteed (ordered tokens, bidirectional scan) &
Universal approximator of permutation-equivariant functions given PE; with
suitable PE/SE more expressive than any WL test, matching graph Transformers;
unbounded power without PE/SE when a random-walk neighborhood encoder is used &
Part of the power is attributed by the authors to neighborhood sampling and
encoding rather than to the selective SSM \\

\midrule

CCMamba~\cite{chen2026ccmamba} &
Rank-aware sequences over incidence neighborhoods &
Upper-bounded by the 1-dimensional combinatorial-complex WL (1-CCWL) test &
An upper bound only: the selective SSM does not exceed higher-order message
passing \\

\midrule

Graph sequence models~\cite{behrouz2024best} &
Depends on tokenizer and backbone &
Task-based analysis beyond WL: recurrent backbones can count node colors when
the state width is at least the number of colors, which non-causal Transformers
without PE cannot; worst-case connectivity favors Transformers (logarithmic
depth), whereas recurrent models need polynomial depth or width unless the
ordering is locality-preserving &
Sensitivity of SSM outputs decays with token distance and deep stacks exhibit
representational collapse; no single backbone dominates \\

\bottomrule
\end{tabular}}
\end{table*}

Three observations follow. First, the available guarantees attach to
components other than the selective scan: GSSC~\cite{huang2024can} obtains its
position between 1-WL and 3-WL from a distance-based kernel, the universality
of GMN~\cite{behrouz2024graph} requires positional encodings exactly as for
graph Transformers, and CCMamba~\cite{chen2026ccmamba} is bounded by the WL
variant of the structure on which it operates. The one result that isolates
selectivity is that the selection mechanism raises the counting power of GSSC
from 3-cycles to 4-cycles. Second, for serialization-based models a comparison
with the WL hierarchy must be interpreted with care. An order-sensitive model
can separate WL-indistinguishable graphs simply because the imposed ordering
breaks symmetries, as random node features do; the same mechanism allows two
isomorphic inputs to receive different outputs, so that distinguishing power
is gained at the expense of invariance. Expressive power should therefore be
characterized under an explicit invariance requirement, for example for a
canonical ordering, in expectation over $\Pi(G)$, or for a multi-order
aggregate. Third, relative to graph Transformers the principal loss is the
fixed-size state. Attention gives every node direct access to every other
node, whereas a selective SSM compresses the prefix of the sequence into a
state of fixed dimension, which limits copying and retrieval from
context~\cite{jelassi2024repeat}; SSMs, like Transformers, are moreover
confined to the complexity class $\mathsf{TC}^0$ \cite{merrill2024illusion}.
The analysis of graph sequence models~\cite{behrouz2024best} indicates that
the comparison is task-dependent, and that the sensitivity of an SSM output to
an input token decays with their distance in the sequence, so that the quality
of the serialization determines which dependencies can be represented in
practice. Open questions include the expressive power of serialization-based
models under an invariance constraint, the state dimension and depth required
for standard graph tasks, and the separate contributions of selectivity,
positional encodings, and local message passing in hybrid GNN--SSM layers.

\subsection{Dynamic, Heterogeneous, and Higher-Order Graphs}
\label{sec:ch-complex-graphs}

% [REV-R21] Signpost added (red): design question addressed by this subsection.
 
Beyond static homogeneous graphs, a single fixed serialization is rarely
sufficient. In dynamic and spatio-temporal graphs, an ordering that is
appropriate at one step may become outdated as the topology evolves;
heterogeneous graphs require type-aware representations; and higher-order
graphs require rank-aware or incidence-based representations. Existing
approaches address these needs through temporal
ordering~\cite{choi2024spot,li2024stg}, dynamic graph
learning~\cite{li2024state,yuan2025dg}, type-aware
tokenization~\cite{pan2024hetegraph}, and rank-aware
serialization~\cite{montagna2024topological,chen2026ccmamba}. These mechanisms
increase representation complexity, and jointly preserving temporal, semantic,
and higher-order structure within efficient selective state-space processing
remains difficult, particularly when these properties coexist.

\subsection{Interpretability and Explainability}
\label{sec:ch-interpretability}

% [REV-R21] Signpost added (red): design question addressed by this subsection.
 
Interpretability is particularly challenging in Graph Mamba because predictions
depend on both graph structure and selective state-space processing. Existing
GNN explanation methods can identify influential nodes, edges, subgraphs, or
features, but they do not directly explain how graph serialization and latent
state evolution affect the final prediction.

Graph Mamba therefore introduces a multi-level explanation problem: an
explanation may need to account for structural importance, sensitivity to the
selected serialization, and information propagation through the SSM. This is
especially important in hybrid architectures, where local graph processing and
global state-space modeling contribute complementary information.

A concrete starting point is the retention profile
$\alpha_{t\leftarrow k}$ of Equation~\eqref{eq:alpha-retention}. Because it
quantifies the contribution of position $k$ to the output at position $t$, it
is itself an explanation primitive and plays for Graph Mamba the role that
attention maps play for Transformers~\cite{ali2025hidden}. Mapping positions
back to nodes through the ordering $\pi$ yields a node-to-node relevance map,
$R_{\pi(t),\pi(k)}=\frac{1}{D}\sum_{d}|\alpha^{(d)}_{t\leftarrow k}|$ for
$k\le t$, which requires no additional training; for bidirectional or
multi-order models the maps of the individual scans are combined. Such maps
cover only the state-space path, not the gating or local message-passing
branches, and they depend on the serialization. They should therefore be
checked for stability across orderings and validated with perturbation tests,
as is required for attention weights~\cite{jain2019attention}. Existing models pursue explanation by other means. One example is FuzzMamba, which introduces fuzzy role memberships
and rule-based reasoning to expose how node behavior and environmental factors
influence spatio-temporal predictions \cite{chen2026modeling}. ExPath provides a more explicit example through PathExplainer, which learns
pathway-level masks to identify the biological subgraphs most responsible for
the model prediction, while PathMamba captures long-range dependencies along
sampled molecular pathways \cite{kotoge2026expath}.

\subsection{Pre-training, Transfer, and Graph Foundation Models}
\label{sec:ch-pretraining}

% [REV-R21] Signpost added (red): design question addressed by this subsection.
Transfer across graphs depends above all on tokenization and ordering, since a pre-trained state-space model encodes the statistics of the sequences it was trained on.
Pre-training followed by adaptation is an established paradigm for GNNs,
through attribute- and structure-level pretext tasks~\cite{hu2019strategies},
cross-network contrastive learning~\cite{qiu2020gcc}, and masked graph
autoencoding~\cite{hou2022graphmae}, and it underlies current work on graph
foundation models~\cite{liu2025graph,mao2024position}. In the Graph Mamba
literature, by contrast, self-supervision is used almost exclusively within a
single dataset: the objectives of DG-Mamba~\cite{yuan2025dg},
Mamba-GTC~\cite{Meng2026Mamba-GTC}, and GLADMamba~\cite{fu2025gladmamba}
% [REV-R23] Undefined reference subsec:learning-strategies replaced by sec:taxonomy (red).
(Section~\ref{sec:taxonomy}) do not produce an encoder that is
reused on other graphs. A domain-specific exception is
Transfer-Mamba~\cite{cheng2025transfer}, which pre-trains a Mamba and
adaptive-GCN encoder by masked reconstruction on traffic data from several
source cities and transfers it to a data-scarce target city.

Several obstacles are specific to the state-space formulation. A pre-trained
selective SSM encodes regularities of the sequences on which it was trained;
if orderings derive from dataset-specific statistics such as degree ranks or
sensor layout, the sequence distribution changes across graphs, so transfer
requires serializers with comparable output statistics or tokenizations based
on recurring local substructures. Masked reconstruction is well matched to
bidirectional attention but not to a causal recurrence, in which a masked
token can only be predicted from its prefix; MambaMIM addresses this in vision
with a token-interpolation scheme~\cite{tang2025mambamim}, and graph-specific
counterparts remain to be developed. Node-feature spaces and spectral
positional encodings are not directly comparable across graphs, and a fixed
state dimension must serve graphs whose sizes differ by orders of magnitude. In domains where GNN pre-training is
already effective, such as brain network analysis~\cite{yang2023ptgb}, a direct
comparison of pre-trained GNN, graph Transformer, and Graph Mamba encoders
would be particularly informative.

\subsection{Hybrid Attention--SSM Designs and Mamba-2/SSD Backbones}
\label{sec:ch-hybrid}

% [REV-R21] Signpost added (red): design question addressed by this subsection.
Backbone choices concern coupling and integration: how the state-space operator is parameterized and how it is combined with attention. Implementation efficiency raises challenges that are specific to
state-space models on graphs. The selective recurrence is evaluated through a
parallel scan rather than a single matrix multiplication, while graph-specific
designs can multiply this cost through bidirectional or multi-order scans.
Moreover, graphs of different sizes produce variable-length token sequences,
making efficient batching dependent on padding or packed scans with state resets
at graph boundaries~\cite{xu2024packmamba}. Graph construction, serialization, and
spectral preprocessing may further dominate the end-to-end cost on large
graphs. These considerations motivate graph-aware scan kernels, packed
variable-length processing, fused multi-direction scans, and sparse or
partitioned serialization.

\paragraph{Hybrid attention--SSM architectures.}
In language modeling, interleaving a few attention layers with Mamba layers has
proved more effective than either component alone:
Jamba~\cite{lieber2024jamba} combines both in a single stack, and a controlled
comparison at the 8B-parameter scale found that a hybrid with a small fraction
of attention layers outperforms a pure Transformer, while pure SSMs lag on
copying and in-context retrieval~\cite{waleffe2024empirical}.
% [REV-R23] Undefined reference sec4 replaced by subsec:static-models (red).
The hybrids
reviewed in Section~\ref{subsec:static-models} are of a different kind: they pair local message
passing with a global SSM, typically by replacing the attention module of a
GraphGPS-style layer~\cite{rampavsek2022recipe}, and thus substitute attention
rather than complement it. Attention is permutation-equivariant and gives
direct access to arbitrary node pairs, which compensates for the fixed-size
state and distance-decaying sensitivity of the SSM, while the SSM supplies an
ordering-based inductive bias and linear-time propagation.
GSM++~\cite{behrouz2024best} instantiates this idea by applying Mamba layers
followed by a Transformer block to hierarchically tokenized graphs and reports
improvements over both pure backbones on most of its benchmarks, and
MapsTSF couples a Mamba-2 backbone with Transformer attention over road-network
traversals~\cite{wang2025mapstsf}. Open design questions include the ratio and
placement of attention layers, whether attention should act on all nodes, on a
sparse pattern~\cite{shirzad2023exphormer}, or on coarsened super-nodes, and
the extent to which a few equivariant attention layers reduce the order
sensitivity $\mathcal{S}(G)$ of Equation~\eqref{eq:order-sensitivity}.

\paragraph{Mamba-2 and SSD for graphs.}
Most surveyed architectures build on the original selective scan. Mamba-2
% [REV-R23] Undefined reference eq:retention-profile replaced by eq:alpha-retention (red).
restricts the state transition to a scalar multiple of the identity,
$\bar{A}_t=a_t\mathbf{I}$, under which the sequence map of
Equation~\eqref{eq:alpha-retention} takes the matrix form

\begin{equation}
\label{eq:ssd-graph}
\mathbf{Y}
=
\left(\mathbf{L}\odot\mathbf{C}\mathbf{B}^{\top}\right)\mathbf{U},
\qquad
L_{tk}=\prod_{j=k+1}^{t}a_j \;\; (t\ge k),
\qquad
L_{tk}=0 \;\; (t<k),
\end{equation}

\noindent where the rows of $\mathbf{C}$ and $\mathbf{B}$ collect $C_t$ and
$\bar{B}_t$, $\mathbf{U}$ stacks the inputs, and $\mathbf{L}$ is a
1-semiseparable mask~\cite{dao2024transformers}. The retention coefficients of
Equation~\eqref{eq:alpha-retention} are then the entries of a masked attention
matrix, and $\mathbf{L}$ makes explicit where the sequential ordering enters
the computation. For graphs, this suggests that SSD offers a faster backbone
with larger state dimensions and support for packed variable-length inputs;
that the sequence-distance decay encoded by $\mathbf{L}$ could be replaced by a
graph-aware structured mask, for instance derived from graph distances, as
long as the mask admits sub-quadratic multiplication (the factorizable kernels
of GSSC~\cite{huang2024can} are close to this idea); that non-causal SSD
variants proposed for images~\cite{shi2025vssd} are natural candidates for
unordered node sets; and that the multi-head structure could be aligned with
node or relation types.
% [REV-R21] Wording corrected (red): Mamba-2 keeps one decay per head; it removes per-state-dimension decay rates, not selectivity.
% [REV-R21-ORIGINAL] Conversely, the scalar-identity transition removes per-channel selectivity, and
Conversely, the scalar-identity transition removes per-state-dimension decay rates, retaining one decay rate per head, and a comparison of sequence backbones including
Mamba and Mamba-2 within one graph pipeline found none to be uniformly
best~\cite{behrouz2024best}; controlled comparisons under identical
serialization and encoding choices are needed.

\subsection{Future Research Directions}
\label{sec7}

Table~\ref{tab:future-directions} summarizes, for each challenge discussed in
this section, promising research directions and representative early efforts.
% [REV-R21] Replaced: three priorities here vs. four directions in the conclusion; both now list the same four priorities, each tied to its design question.
% [REV-R21-ORIGINAL] Three priorities emerge. The first is \emph{reliability}: order sensitivity,
% [REV-R21-ORIGINAL] robustness, and behavior under distribution shift should be measured routinely
% [REV-R21-ORIGINAL] and, where possible, controlled by design. The second is \emph{theory}: the
% [REV-R21-ORIGINAL] expressive power of serialization-based models should be characterized under
% [REV-R21-ORIGINAL] an invariance constraint, so that Graph Mamba can be positioned rigorously
% [REV-R21-ORIGINAL] relative to the WL hierarchy and to graph Transformers. The third is
% [REV-R21-ORIGINAL] \emph{infrastructure and scale}: graph-aware scan primitives, SSD-based
% [REV-R21-ORIGINAL] backbones, hybrid attention--SSM layers, and transferable pre-training are
% [REV-R21-ORIGINAL] jointly required before SSM-based encoders can serve as graph foundation
% [REV-R21-ORIGINAL] models.
Four priorities emerge, one for each group of design questions, and they structure the conclusion. The first is \emph{order-robust modeling} (ordering): order sensitivity, robustness, and behavior under distribution shift should be measured routinely and, where possible, controlled by design. The second is \emph{theory} (tokenization and coupling): the expressive power of serialization-based models should be characterized under an explicit invariance requirement, relative to the WL hierarchy and to graph Transformers. The third is \emph{infrastructure and backbone design} (coupling and integration): graph-aware scan primitives, Mamba-2/SSD backbones with graph-aware structured masks, and hybrid attention--SSM layers. The fourth is \emph{transferable pre-training} (tokenization and ordering): serializers and tokenizers whose output statistics carry over across graphs, so that SSM-based encoders can serve as graph foundation models.

\begin{table*}[!htbp]
\centering
\caption{Open challenges and future research directions for Graph Mamba.}
\label{tab:future-directions}
\small
\renewcommand{\arraystretch}{1.12}
\resizebox{\textwidth}{!}{%
% [REV-R21] New column "Design question" (red); widths adjusted.
% [REV-R21-ORIGINAL] \begin{tabular}{>{\raggedright\arraybackslash}p{3.0cm} >{\raggedright\arraybackslash}p{9.2cm} >{\raggedright\arraybackslash}p{5.2cm}}
\begin{tabular}{>{\raggedright\arraybackslash}p{3.0cm} >{\raggedright\arraybackslash}p{2.4cm} >{\raggedright\arraybackslash}p{7.4cm} >{\raggedright\arraybackslash}p{4.6cm}}
\toprule
\textbf{Challenge} & \textbf{Design question} & \textbf{Promising directions} & \textbf{Early work / examples} \\
\midrule
End-to-end scalability and tooling & Ordering, integration &
Sparse graph-aware scanning, efficient serialization, graph partitioning;
packed variable-length scans with state resets at graph boundaries, fused
bidirectional and multi-order scans, portable reference implementations &
GSSC~\cite{huang2024can}, GraphSSM~\cite{li2024state};
PackMamba~\cite{xu2024packmamba} and SSD kernels~\cite{dao2024transformers}
(sequence domain) \\
\midrule
Serialization and permutation & Ordering, coupling &
Learnable ordering, permutation-equivariant SSM operators, multi-order scans,
topology-conditioned state transitions &
Graph-Mamba~\cite{wang2024graph}, GMN~\cite{behrouz2024graph},
GSSC~\cite{huang2024can}, GrassNet~\cite{zhao2024grassnet} \\
\midrule
Robustness and distribution shift & Ordering &
Serializer-aware attacks and certified defenses, order-sensitivity metrics
(Equation~\eqref{eq:order-sensitivity}), order-consistency regularization,
evaluation on graph out-of-distribution (OOD) benchmarks &
DG-Mamba~\cite{yuan2025dg}; spectral-norm constraints in
DyG-Mamba~\cite{li2024dyg} \\
\midrule
Expressive power & Tokenization, coupling &
Expressivity under canonical or expected orderings, state-dimension and depth
bounds for graph tasks, isolating the roles of selectivity and of PE/SE &
GSSC~\cite{huang2024can}, GMN~\cite{behrouz2024graph},
CCMamba~\cite{chen2026ccmamba}, graph sequence models~\cite{behrouz2024best} \\
\midrule
Dynamic and complex graphs & Tokenization, ordering &
Adaptive state transitions, type-aware state updates, temporal structural
conditioning, rank-aware representations &
GraphSSM~\cite{li2024state}, DyGMamba~\cite{ding2024dygmamba},
HeteGraph-Mamba~\cite{pan2024hetegraph},
TopoMamba~\cite{montagna2024topological}, CCMamba~\cite{chen2026ccmamba} \\
\midrule
Interpretability & Coupling, integration &
Retention-based node relevance maps validated by perturbation tests, stability
of explanations across orderings, component-level explanations covering the
local and gating branches &
Hidden-attention view of Mamba~\cite{ali2025hidden};
FuzzMamba~\cite{chen2026modeling}; ExPath~\cite{kotoge2026expath} \\
\midrule
Self-supervision, pre-training, and foundation models & Tokenization, ordering &
Graph-specific pretext tasks suited to causal and bidirectional scans,
transferable serializers and substructure tokenizers, cross-graph comparable
encodings, size and length generalization, scaling studies &
DG-Mamba~\cite{yuan2025dg}, Mamba-GTC~\cite{Meng2026Mamba-GTC},
GLADMamba~\cite{fu2025gladmamba}; Transfer-Mamba~\cite{cheng2025transfer} \\
\midrule
Hybrid attention--SSM and SSD backbones & Coupling, integration &
Sparse or coarsened attention interleaved with SSM layers, graph-aware
structured masks in the SSD form (Equation~\eqref{eq:ssd-graph}), non-causal
SSD for unordered nodes, controlled Mamba vs.\ Mamba-2 comparisons &
GSM++~\cite{behrouz2024best}, MapsTSF~\cite{wang2025mapstsf};
Jamba~\cite{lieber2024jamba} and VSSD~\cite{shi2025vssd} outside graphs \\
\midrule
Cross-modal graph modeling & Integration &
Cross-modal state conditioning, graph-token alignment, multimodal selective
fusion &
M3-HGMA~\cite{gubbala2026m3}, BrainMamba~\cite{behrouz2024brain} \\
\bottomrule
\end{tabular}}
\end{table*}


 

\section{Conclusion}
\label{sec8}

 This survey examined how Mamba and selective state-space models are adapted to
graph-structured data. The reviewed architectures differ less in the
state-space component itself than in how they address four core design
questions: what the model reads through tokenization, in what order it processes
graph elements, how strongly graph structure conditions the state dynamics, and
how state-space processing is integrated with graph propagation. Static graphs
mainly emphasize ordering, dynamic graphs emphasize structural conditioning,
and heterogeneous or higher-order graphs place greater demands on tokenization.
Overall, most existing designs still incorporate graph information primarily
before or around the state-space block rather than directly within its dynamics.

These design choices have important consequences. Serialization introduces
sensitivity to node ordering and complicates permutation consistency, while the
linear complexity of selective scanning does not automatically translate into
linear end-to-end complexity once graph construction, tokenization, structural
encoding, and fusion are considered. Published studies nevertheless report
competitive performance across healthcare, remote sensing, spatio-temporal
forecasting, finance, language, and other graph-learning tasks. However, the
current evidence remains heterogeneous, and improvements cannot yet be
attributed unambiguously to ordering strategies, graph--SSM coupling, or the
state-space component itself. Thus, Graph Mamba has demonstrated promising
efficiency and broad applicability, while its advantages for long-range graph
reasoning still require more controlled evaluation.

Future progress will depend on improving order robustness and permutation
consistency, establishing a clearer theory of expressive power under invariance
constraints, developing graph-aware Mamba-2/SSD and hybrid attention--SSM
backbones, and enabling transferable pre-training across graphs and domains.
Addressing these challenges will clarify when selective state-space modeling
provides a genuine advantage over conventional GNNs and graph Transformers and
help establish Graph Mamba as a more mature and reliable graph-learning
paradigm.


\section*{Acknowledgement}
The authors gratefully acknowledge Prince Sultan University, Riyadh, Saudi Arabia, for its support of this research.

%==================REFERENCES====================================================
%\bibliographystyle{elsarticle-num} 
\bibliographystyle{ACM-Reference-Format}
\bibliography{cas-refs}

\clearpage

%====================================================================== 
\appendix

\section{Detailed Performance Comparisons}
\label{app:performance-tables}

This appendix provides the detailed numerical comparisons referenced in
Section~\ref{sec:performance-comparison}.

\begin{table}[h]
\centering
\caption{
Mean classification accuracy (\%) on brain benchmarks.
% [REV-R20] Source of the comparison added (red).
Results as reported in~\cite{behrouz2024brain}.
}
\label{tab:Healthcompare}
\scriptsize

\begin{tabular}{lccc}
\toprule
\textbf{Model} &
\textbf{BVFC-MEG} &
\textbf{HCP-Mental} &
\textbf{HCP-Age} \\
\midrule
USAD~\cite{audibert2020usad}
& $50.02\pm1.13$ & $73.49\pm1.56$ & $39.17\pm1.68$ \\
GMM~\cite{he2023generalization}
& $53.04\pm1.73$ & $90.92\pm1.83$ & $47.75\pm1.26$ \\
GraphMixer~\cite{cong2023we}
& $53.12\pm1.18$ & $91.13\pm1.44$ & $48.32\pm1.11$ \\
BrainNetCNN~\cite{kawahara2017brainnetcnn}
& $50.12\pm1.57$ & $83.58\pm1.68$ & $42.26\pm2.03$ \\
BrainGNN~\cite{li2021braingnn}
& $51.08\pm0.96$ & $85.25\pm2.17$ & $43.08\pm1.54$ \\
FBNetGen~\cite{kan2022fbnetgen}
& $50.94\pm1.39$ & $84.47\pm1.88$ & $42.83\pm1.78$ \\
ADMIRE~\cite{behrouz2023admire++}
& $54.87\pm1.92$ & $89.74\pm1.93$ & $47.82\pm1.72$ \\
PTGB~\cite{yang2023ptgb}
& $55.11\pm1.62$ & $92.58\pm1.31$ & $48.41\pm1.47$ \\
BNTransformer~\cite{kan2022brain}
& $55.17\pm1.74$ & $91.71\pm1.48$ & $47.94\pm1.15$ \\
BrainMixer~\cite{behrouz2023learning}
& $62.58\pm1.12$ & $96.32\pm0.29$ & $57.83\pm1.03$ \\
\rowcolor{lightgray}
GraphS4mer~\cite{tang2023modeling}
& $74.39\pm0.92$ & $82.30\pm1.83$ & $46.32\pm1.09$ \\
\rowcolor{lightgray}
BrainMamba~\cite{behrouz2024brain}
& $78.03\pm1.69$ & $96.57\pm1.05$ & $59.62\pm1.71$ \\
\bottomrule
\end{tabular}
\end{table}

\begin{table*}[h]
\centering
\caption{
Published performance comparison on hyperspectral image classification.
% [REV-R20] Source of the comparison added (red).
Results as reported in~\cite{yang2024graphmamba}.
}
\label{tab:RScomparison}
\scriptsize
\begin{tabular}{lcccccccccc}
\toprule
\multirow{2}{*}{\textbf{Model}} &
\multicolumn{3}{c}{\textbf{Indian Pines}} &
\multicolumn{3}{c}{\textbf{Salinas}} &
\multicolumn{3}{c}{\textbf{UH2013}} &
\multirow{2}{*}{\textbf{Complexity (G)}} \\
\cmidrule(lr){2-4}
\cmidrule(lr){5-7}
\cmidrule(lr){8-10}
&
\textbf{OA} & \textbf{AA} & \textbf{Kappa} &
\textbf{OA} & \textbf{AA} & \textbf{Kappa} &
\textbf{OA} & \textbf{AA} & \textbf{Kappa} & \\
\midrule

AB-LSTM~\cite{mei2021hyperspectral}
& 56.68 & 60.41 & 51.02
& 74.83 & 83.89 & 72.21
& 79.61 & 82.38 & 77.97 & -- \\

3D CNN~\cite{hamida2018deep}
& 65.83 & 79.12 & 61.38
& 87.59 & 92.96 & 86.21
& 71.84 & 73.84 & 69.54 & 0.69 \\

GAHT~\cite{mei2022hyperspectral}
& 85.96 & 92.96 & 84.05
& 94.21 & 97.37 & 93.56
& 90.84 & 92.13 & 90.10 & 4.27 \\

\rowcolor{lightgray}
GraphMamba (HSI)~\cite{yang2024graphmamba}
& 96.43 & 97.83 & 95.91
& 97.33 & 98.66 & 97.02
& 95.62 & 96.23 & 95.27 & 0.69 \\

\bottomrule
\end{tabular}
\end{table*}

\begin{table*}[h]
\centering
\caption{
Published performance comparison on spatio-temporal forecasting benchmarks.
% [REV-R20] Source of the comparison added (red).
Results as reported in~\cite{li2024stg}.
}
\label{tab:TRperformance_comparison}
\scriptsize
\begin{tabular}{lccccccccc}
\toprule
\multirow{2}{*}{\textbf{Model}} &
\multicolumn{3}{c}{\textbf{PeMS04}} &
\multicolumn{3}{c}{\textbf{HZMetro}} &
\multicolumn{3}{c}{\textbf{KnowAir}} \\
\cmidrule(lr){2-4}
\cmidrule(lr){5-7}
\cmidrule(lr){8-10}
&
\textbf{RMSE} & \textbf{MAE} & \textbf{MAPE (\%)} &
\textbf{RMSE} & \textbf{MAE} & \textbf{MAPE (\%)} &
\textbf{RMSE} & \textbf{MAE} & \textbf{MAPE (\%)} \\
\midrule

% [REV-NOTE] This row merges three different methods (Seq2Seq, PVCGN, air-quality neural ODE); only the first is an LSTM. The HZMetro MAPE (9.90) is lower than all other rows; check the values.
LSTM-based methods
~\cite{sutskever2014sequence,liu2020physical,tian2024air}
& 40.49 & 26.24 & 19.30
& 57.10 & 35.27 & 9.90
& 30.88 & 20.21 & -- \\

STGCN~\cite{yu2017stgcn}
& 35.55 & 22.70 & 14.59
& 34.85 & 21.33 & 13.47
& 11.46 & 8.37 & 11.26 \\

STAEformer~\cite{liu2023spatio}
& 30.18 & 18.22 & 11.98
& 29.94 & 18.85 & 12.03
& 8.69 & 6.93 & 9.89 \\

\rowcolor{lightgray}
STG-Mamba~\cite{li2024stg}
& 29.53 & 18.09 & 12.11
& 29.23 & 18.26 & 11.59
& 8.05 & 6.37 & 9.64 \\

\bottomrule
\end{tabular}
\end{table*}
\begin{table}[h]
\centering
\caption{
Published performance comparison for financial forecasting.
% [REV-R20] Source of the comparison added (red).
Results as reported in~\cite{mehrabian2024mamba}.
}
\label{tab:FMcomparison}
\scriptsize

\begin{tabular}{lcccc}
\toprule
\textbf{Model} &
\textbf{NASDAQ} &
\textbf{NYSE} &
\textbf{DJIA} &
\textbf{Training time (s/epoch)} \\
\midrule

LSTM~\cite{moghar2020stock}
& 0.0187 & 0.0144 & 0.0121 & 0.081 \\

Transformer~\cite{wang2022stock}
& 0.0147 & 0.0141 & 0.0166 & 1.91 \\

FourierGNN~\cite{yi2024fouriergnn}
& 0.0152 & 0.0146 & 0.0190 & 1.10 \\

\rowcolor{lightgray}
SAMBA~\cite{mehrabian2024mamba}
& 0.0128 & 0.0125 & 0.0108 & 0.891 \\

\bottomrule
\end{tabular}
\vspace{0.5mm}

{\footnotesize RMSE is reported for NASDAQ, NYSE, and DJIA.}
\end{table}

\begin{table*}[h]
\centering
\caption{
Published performance comparison for aspect-based sentiment analysis.
% [REV-R20] Source of the comparison added (red).
Results as reported in~\cite{lawan2024mambaforgcn}.
}
\label{tab:SAperformance}
\scriptsize
\begin{tabular}{lcccccc}
\toprule
\multirow{2}{*}{\textbf{Model}} &
\multicolumn{2}{c}{\textbf{Restaurant14}} &
\multicolumn{2}{c}{\textbf{Laptop14}} &
\multicolumn{2}{c}{\textbf{Twitter14}} \\
\cmidrule(lr){2-3}
\cmidrule(lr){4-5}
\cmidrule(lr){6-7}
&
\textbf{Accuracy} & \textbf{F1} &
\textbf{Accuracy} & \textbf{F1} &
\textbf{Accuracy} & \textbf{F1} \\
\midrule

CNN over BERT-GCN~\cite{phan2022aspect}
& 85.25 & 78.76
& 78.59 & 74.76
& 74.03 & 72.39 \\

Dual-gated GCN~\cite{liu2023enhancing}
& 83.66 & 76.73
& 75.70 & 72.57
& 74.87 & 72.27 \\

DGEDT~\cite{tang2020dependency}
& 83.90 & 75.10
& 76.80 & 72.30
& 74.80 & 73.40 \\

\rowcolor{lightgray}
MambaForGCN~\cite{lawan2024mambaforgcn}
& 86.68 & 80.86
& 81.80 & 78.59
& 77.67 & 76.88 \\

\bottomrule
\end{tabular}
\end{table*}


\end{document}
\endinput
%%
%% End of file `elsarticle-template-num.tex'.
